% ===========================================================================
% The open/closed realization: cochain derived center and physical bulk
% ===========================================================================
%
% The chiral Hochschild complex and the bordered FM compactification
% each produce an open/closed algebraic structure.  The cochain-level
% derived center answers the universal local-action question for a
% boundary chiral algebra.  A physical bulk factorization algebra enters
% only through a typed comparison map to that cochain object.

\section{The open/closed realization: cochain derived center and physical bulk}
\label{sec:thqg-open-closed-realization}
\index{derived center!chiral|textbf}
\index{open/closed realization|textbf}
\index{Swiss-cheese theorem!chiral|textbf}
\index{cochain-level closed-sector actor|textbf}
\index{algebraic sector|textbf}
\index{BTZ hypothesis|textbf}
\index{Cardy hypothesis|textbf}
\index{Maloney--Witten Farey-tail|textbf}

\begin{remark}[Algebraic sector identification;
\ClaimStatusProvedHere]
\label{rem:thqg-oc-algebraic-sector}
\textit{Type signature.}
Quadrant: Open. Presentation: chain-level + $(\infty,1)$.
Level: $1 \to 3$ (chart $\to$ derived chiral centre).
Hypothesis package: complete curved $A_\infty$-chiral algebra
\textup{(}arity and pole-order filtrations complete and
separated\textup{)}; chiral Maurer--Cartan equation $m\{m\} = 0$.
\medskip

The holographic register inscribed in this section is the
\emph{algebraic sector}:
\begin{equation}\label{eq:thqg-oc-algebraic-sector}
\partial \;=\; \cA, \qquad
\mathrm{bulk} \;=\; Z^{\mathrm{der}}_{\mathrm{ch}}(\cA), \qquad
\mathrm{interaction} \;=\; \mathsf{SC}^{\mathrm{ch,top}}\text{-brace}.
\end{equation}
The boundary chart is the chiral algebra, the bulk slot is the
cochain-level chiral derived centre
(Definition~\ref{def:thqg-chiral-derived-center}), and the interaction
is the two-colour
$\mathsf{SC}^{\mathrm{ch,top}}$-brace structure on the derived-centre
pair $(Z^{\mathrm{der}}_{\mathrm{ch}}(\cA),\,\cA)$. The chiral
Swiss-cheese Theorem~\ref{thm:thqg-swiss-cheese} realizes the
universal local mixed action; the corresponding cross-volume
inscription is the Vol.~II
$\mathsf{SC}^{\mathrm{ch,top}}$-brace bulk theorem at level~$3$ of the
Beilinson tower, and the Vol.~III propagation is the two-stage
functor $\Phi_d^{(\Sigma_{d-1}, C)} = \mathrm{Sp}^{\mathrm{ch}}_{\Sigma_{d-1},C}
\circ \Phi_d^{\mathrm{FA}}$ on the CY-quadrant side.

The algebraic sector does not by itself construct a compact
dynamical-metric path integral. Level-5 scalar shadows (modular form,
$\kappa$-tuple, partition function) descend from level-4 brane /
operator data by modular trace and clutching; the reverse direction
is the operator-level Pfaffian / square-root reconstruction problem
and remains open. A gravity-line completion of the algebraic sector
requires the BTZ / Cardy / Maloney--Witten / Manschot--Moore Farey-tail
hypothesis package: modular invariance of the boundary partition
function, vacuum (identity-block) dominance in the high-temperature
Cardy channel~\cite{Cardy86}, saddle stability of the BTZ
instanton~\cite{BTZ92}, and the Maloney--Witten / Manschot--Moore
Farey-tail completion~\cite{Maloney-Witten07} of the
$\mathrm{SL}(2,\mathbb{Z})$-orbit sum. With these four inputs supplied,
the Vol.~II gravity-line completion target is the Virasoro-vacuum
scalar shadow $Z_{\mathrm{3d\,QG,\,Vir}}(S^2 \times S^1_\tau) =
q^{(1-c)/24}\,(1-q)/\eta(\tau)$ attached to
$Z^{\mathrm{der}}_{\mathrm{ch}}(\mathrm{Vir}_c)$.
\end{remark}

\begin{remark}[Local input]
\label{rem:thqg-oc-dependency-spine}
The local open/closed construction uses:
\begin{enumerate}[label=\textup{(\roman*)}]
\item the chiral Hochschild complex
 $C^\bullet_{\mathrm{chiral}}(\cA)$
 (Definition~\ref{def:chiral-hochschild}),
\item the bordered FM compactification
 $\overline{C}{}^{\,\mathrm{oc}}_{k,\vec{m}}(X, D, \tau)$
 (\S\ref{sec:bordered-fm}),
\item the open-closed convolution algebra
 $\mathfrak{g}^{\mathrm{oc}}_{\cA, \vec{\mathcal{M}}}$
 (Definition~\ref{def:oc-modular-convolution}), and
\item the bar-intrinsic MC element $\Theta_\cA$
 (Theorem~\ref{thm:mc2-bar-intrinsic}).
\end{enumerate}
The chiral Hochschild complex and bordered FM compactification
govern the genus-$0$ local theory. The gravitational input of
Definition~\ref{def:thqg-standing-hypotheses} enters only when
the modular curvature and clutching terms are invoked.
\end{remark}

% ===================================================================
\subsection{The chiral endomorphism operad and the algebraic derived center}
\label{subsec:thqg-chiral-endomorphism-operad}
\index{chiral endomorphism operad|textbf}
% ===================================================================

The geometric chiral Hochschild complex of
Definition~\ref{def:chiral-hochschild} lives on FM
compactifications with log forms. Its algebraic
counterpart, constructed below, makes the operadic braces and the universal
open/closed structure explicit.

\begin{definition}[Chiral endomorphism operad]
\label{def:thqg-chiral-endomorphism-operad}
Let $A$ be a graded vector space equipped with a translation operator
$\partial \colon A \to A$ of degree~$0$.
The \emph{chiral endomorphism operad} $\mathcal{E}\!nd^{\mathrm{ch}}_A$
is the nonsymmetric operad with:
\begin{enumerate}[label=\textup{(\alph*)}]
\item \emph{Components.}
 $\mathcal{E}\!nd^{\mathrm{ch}}_A(n)
 := \operatorname{Hom}\bigl(A^{\otimes n},\,
 A((\lambda_1))\cdots((\lambda_{n-1}))\bigr)$
 for $n \ge 1$, where $\lambda_1, \ldots, \lambda_{n-1}$ are
 formal spectral variables of degree~$0$.

\item \emph{Operadic composition.} For $f \in \mathcal{E}\!nd^{\mathrm{ch}}_A(n)$
 and $g \in \mathcal{E}\!nd^{\mathrm{ch}}_A(m)$, the $i$-th partial
 composition $f \circ_i g \in \mathcal{E}\!nd^{\mathrm{ch}}_A(n + m - 1)$
 is defined by
 \begin{equation}\label{eq:thqg-block-composition}
 (f \circ_i g)(a_1, \ldots, a_{n+m-1})
 :=
 f(a_1, \ldots, a_{i-1},\,
 g(a_i, \ldots, a_{i+m-1}),\,
 a_{i+m}, \ldots, a_{n+m-1})
 \end{equation}
 with the \emph{block-substitution rule} for spectral variables:
 the variable $\lambda_i$ in $f$ is replaced by the block sum
 $\Lambda_i := \lambda_i + \lambda_{i+1} + \cdots + \lambda_{i+m-2}$,
 and the variables $\lambda_{i}, \ldots, \lambda_{i+m-2}$ in $g$ are
 inserted consecutively.
\end{enumerate}
Associativity of the operadic composition follows from associativity
of consecutive block-substitution: if a block $B_1$ is inserted into
a block $B_2$ which is itself inserted into an ambient expression,
the result is a single contiguous block $B_1 \subset B_2$ whose
variables partition correctly.
\end{definition}

\begin{remark}[Relation to the OPE]
\label{rem:thqg-ope-relation}
For a vertex algebra $V$ with OPE
$a(z)\, b(w) \sim \sum_{n \ge 0} (a_{(n)} b)(w)\, (z - w)^{-n-1}$,
the element $\mu \in \mathcal{E}\!nd^{\mathrm{ch}}_V(2)$ encoding
the OPE is
\[
\mu(a, b; \lambda)
=
\sum_{n \ge 0} (a_{(n)} b)\, \lambda^n / n!
=
a_{(\lambda)} b,
\]
which is exactly the $\lambda$-bracket of the corresponding
Lie conformal algebra. The higher operations
$\mu_n \in \mathcal{E}\!nd^{\mathrm{ch}}_V(n)$ for $n \ge 3$ encode
iterated OPE with the spectral variables tracking the collision
parameters.
\end{remark}

\begin{proposition}[Bridge: BD chiral operad $\leftrightarrow$ algebraic $\mathcal{E}\!\mathit{nd}^{\mathrm{ch}}$]
\label{prop:bd-algebraic-bridge}
\ClaimStatusProvedHere
Let $\cA$ be a chiral algebra on a smooth curve~$X$, and let
$p \in X$ with local coordinate $z \colon \widehat{D}_p
\xrightarrow{\sim} \operatorname{Spec} \mathbb{C}[[z]]$. Assume
that the completed finite-piece of the underlying right
$\mathcal D$-module of~$\cA$ is locally free with flat connection
at~$p$. The horizontal fibre
$A_p := \Gamma_{\nabla}(\widehat{D}_p,\cA)$ carries the translation
operator $\partial=\nabla_{d/dz}$. The construction below produces
a coordinate-dependent isomorphism of non-$\Sigma$ operads:
\begin{equation}\label{eq:bd-algebraic-bridge}
 \operatorname{End}^{\mathrm{ch}}_{\cA, \mathrm{BD}}(n)\big|_{\widehat{D}_p^n}
 \;\xrightarrow[\;\cong\;]{\;\Phi_z\;}
 \mathcal{E}\!\mathit{nd}^{\mathrm{ch}}_{A_p}(n)
\end{equation}
where the left side is the Beilinson--Drinfeld chiral
endomorphism operad
\textup{(Definition~\textup{\ref{def:chiral-endo}})} restricted
to the formal polydisk, and the right side is the algebraic
chiral endomorphism operad
\textup{(Definition~\textup{\ref{def:thqg-chiral-endomorphism-operad}})}.

The isomorphism $\Phi_z$ is constructed in four steps:
\begin{enumerate}[label=\textup{Step~\arabic*.},leftmargin=*]
\item \textbf{Formal-disk restriction.}
 Restrict the BD operad from $X^n$ to the formal
 polydisk $\widehat{D}_p^n = \operatorname{Spec} \mathbb{C}[[z_1, \ldots, z_n]]$.
\item \textbf{D-module trivialisation.}
 The chosen coordinate identifies the completed locally free
 flat $\mathcal D$-module with
 $A_p\widehat{\otimes}\mathbb{C}[[z]]$. This is a trivialisation
 of the flat finite-piece attached to~$\cA$, not a statement that
 arbitrary coherent $\mathcal D$-modules on a formal disk are
 trivial. The chiral algebra $\cA|_{\widehat{D}_p}$ becomes the
 graded vector space $A_p$ equipped with
 $\partial = \nabla_{d/dz}$.
\item \textbf{Spectral-variable identification.}
 Under trivialisation, the maximal meromorphic extension
 $j_* j^*(\cA^{\boxtimes n})|_{\widehat{D}_p^n}$ becomes
 $A_p^{\otimes n} \otimes
 \mathbb{C}((\lambda_1))\cdots((\lambda_{n-1}))$
 where $\lambda_i = z_i - z_n$ are the relative positions.
 The BD chiral product $\mu^{\mathrm{ch}}$ restricts to the
 $\lambda$-bracket: $\mu(a, b; \lambda) = a_{(\lambda)} b$.
\item \textbf{Operadic composition.}
 The operadic composition on the BD side (via boundary strata
 of the FM compactification $\overline{C}_n(\widehat{D}_p)$) maps
 to block-substitution on the algebraic side: nested collision
 of $m$ points near the $i$-th point of an $n$-point
 configuration corresponds to the spectral substitution rule
 $\lambda_i \mapsto \Lambda_i = \lambda_i + \cdots +
 \lambda_{i+m-2}$.
\end{enumerate}
The isomorphism $\Phi_z$ depends on the choice of local
coordinate~$z$. The group $\operatorname{Aut}(\mathcal{O})$ of formal coordinate
changes acts on $\mathcal{E}\!\mathit{nd}^{\mathrm{ch}}_{A_p}(n)$ and
provides descent data: the BD operad is the
$\operatorname{Aut}(\mathcal{O})$-equivariant global object from which the algebraic
operad is obtained by choosing a trivialisation.
\end{proposition}

\begin{proof}
Step~2 is the standard formal-disk trivialisation of a locally
free module with flat connection: a section is identified with its
horizontal value at~$p$, and the coordinate vector field gives the
translation operator. Step~3 is the definition of the
$j_* j^*$-functor applied to the trivialised finite-piece: sections
with poles along the diagonal $z_i = z_j$ become completed formal
Laurent series in the difference variables $\lambda_i = z_i - z_n$.
Poles along non-last diagonals $z_i = z_j$ ($i, j < n$) are
captured by the element $(\lambda_i - \lambda_j)^{-1}$ in the
iterated Laurent field; the ordering of the iterated completion
(expand in $\lambda_{n-1}$ first, then $\lambda_{n-2}$, etc.)
encodes the collision ordering, and this is the source of the
non-$\Sigma$ (ordered) operadic structure.
Step~4 follows from the operadic structure of the
Fulton--MacPherson compactification~\cite{FM94}: the formal-disk
restriction identifies codimension-one boundary strata with
consecutive-block collisions, matching the algebraic operadic
composition in equation~\eqref{eq:thqg-block-composition}. The
$\operatorname{Aut}(\mathcal{O})$-equivariance follows from the
coordinate naturality of each step: a formal coordinate change
$z \mapsto \tilde{z}(z)$ transforms $\partial$, the spectral
variables, and the block-substitution rule compatibly.
\end{proof}

\begin{definition}[Chiral Hochschild cochain complex, algebraic model]
\label{def:thqg-algebraic-hochschild}
\index{chiral Hochschild cochains!algebraic model|textbf}
Let $\cA$ be a chiral algebra with underlying graded vector space $A$
and $A_\infty$-chiral structure
$m = \sum_{n \ge 1} m_n \in \prod_{n \ge 1}
\mathcal{E}\!nd^{\mathrm{ch}}_A(n)$.
The \emph{algebraic chiral Hochschild cochain complex} is
\begin{equation}\label{eq:thqg-algebraic-hochschild}
\mathcal{C}^{\bullet}_{\mathrm{ch}}(\cA, \cA)
\;:=\;
\prod_{n \ge 0}
\mathcal{E}\!nd^{\mathrm{ch}}_A(n+1)[-n],
\end{equation}
so a degree-$p$ element $f$ is a sequence
$(f_0, f_1, f_2, \ldots)$ with
$f_n \in \mathcal{E}\!nd^{\mathrm{ch}}_A(n+1)$ of internal degree
$p + n$.
\end{definition}

\begin{definition}[Operadic braces]
\label{def:thqg-operadic-braces}
\index{brace algebra!chiral|textbf}
For $f \in \mathcal{E}\!nd^{\mathrm{ch}}_A(n+1)$ and
$g_1 \in \mathcal{E}\!nd^{\mathrm{ch}}_A(k_1 + 1), \ldots,
g_r \in \mathcal{E}\!nd^{\mathrm{ch}}_A(k_r + 1)$,
define the \emph{brace operation}
\begin{equation}\label{eq:thqg-brace}
f\{g_1, \ldots, g_r\}
\;:=\;
\sum_{1 \le i_1 < \cdots < i_r \le n}
(-1)^{\varepsilon}
f \circ_{i_1} g_1 \circ_{i_2 + k_1} g_2
\circ_{i_3 + k_1 + k_2} g_3
\cdots \circ_{i_r + k_1 + \cdots + k_{r-1}} g_r,
\end{equation}
where $\varepsilon = \sum_{s=1}^{r}
(|g_s| - 1)(i_s + k_1 + \cdots + k_{s-1} - s)$ is the
Koszul sign from permuting $g_s$ past the preceding inputs.
The one-fold brace defines the \emph{Gerstenhaber bracket}:
\begin{equation}\label{eq:thqg-gerstenhaber}
[f, g]
\;:=\;
f\{g\} - (-1)^{(|f|-1)(|g|-1)} g\{f\}.
\end{equation}
The differential on
$\mathcal{C}^{\bullet}_{\mathrm{ch}}(\cA, \cA)$ is
$\delta f := [m, f]$, where
$m = \sum_{n \ge 1} m_n$ is the $A_\infty$-chiral
structure.
\end{definition}

\begin{theorem}[Brace dg algebra structure on chiral Hochschild cochains;
\ClaimStatusProvedHere]
\label{thm:thqg-brace-dg-algebra}
\index{brace algebra!chiral!theorem|textbf}
Let $\cA$ be a chiral algebra. The complex
$(\mathcal{C}^{\bullet}_{\mathrm{ch}}(\cA, \cA), \delta)$
with braces~\eqref{eq:thqg-brace} is a brace dg algebra.
Equivalently:
\begin{enumerate}[label=\textup{(\roman*)}]
\item $\delta^2 = 0$ (using the $A_\infty$-chiral
 relation $m\{m\} = 0$).

\item The \emph{higher pre-Lie identity}:
 for all $f, g_1, \ldots, g_r, h$,
 \begin{equation}\label{eq:thqg-higher-pre-lie}
 f\{g_1, \ldots, g_r\}\{h\}
 =
 \sum_{s=0}^{r}
 (-1)^{(|h|-1)\sum_{t>s}(|g_t|-1)}\,
 f\{g_1, \ldots, g_s, h, g_{s+1}, \ldots, g_r\}
 + \text{insertion terms}.
 \end{equation}

\item The \emph{Leibniz rule}:
 $\delta(f\{g_1, \ldots, g_r\})
 = (\delta f)\{g_1, \ldots, g_r\}
 + \sum_{s} (-1)^{\varepsilon_s}
 f\{g_1, \ldots, \delta g_s, \ldots, g_r\}
 + \text{boundary terms}$.

\item The \emph{Gerstenhaber algebra on cohomology}:
 $H^*(\mathcal{C}^{\bullet}_{\mathrm{ch}}(\cA, \cA), \delta)$
 carries a Gerstenhaber algebra structure, with cup product
 from braces and Lie bracket from~\eqref{eq:thqg-gerstenhaber}.
\end{enumerate}
\end{theorem}

\begin{proof}
\emph{Step~1: $\delta^2 = 0$.}
By direct computation, $\delta^2(f) = [m, [m, f]]$.
In the shifted complex
\eqref{eq:thqg-algebraic-hochschild}, each
$m_n \in \mathcal{E}\!nd^{\mathrm{ch}}_A(n)$ has internal
degree $2-n$, hence cochain degree~$1$ in
$\mathcal{E}\!nd^{\mathrm{ch}}_A(n)[1-n]$. The suspended
degree entering the Gerstenhaber sign is therefore
$|m|-1 = 0$. Thus the graded-Jacobi shortcut
$[m, [m, f]] = \tfrac{1}{2}[[m, m], f]$ is not available
here: at even suspended degree the Jacobi identity only gives
the tautology $[[m, m], f] = 0$.

Instead one uses the brace form of the $A_\infty$ relations.
Since $|m|-1 = 0$, the bracket is
$[m, g] = m\{g\} - g\{m\}$, so
\begin{align*}
\delta^2(f)
&= m\{m\{f\} - f\{m\}\} - (m\{f\} - f\{m\})\{m\} \\
&= m\{m\{f\}\} - m\{f\{m\}\}
 - m\{f\}\{m\} + f\{m\}\{m\}.
\end{align*}
Expanding these four terms by direct insertions, every summand is
either a term containing $m\{m\}$ or a pair of nested/external
insertions with opposite signs. The $A_\infty$-chiral relations
are precisely
\[
(m\{m\})_n
=
\sum_{\substack{r+s+t=n \\ s \ge 1}}
(-1)^{rs+t}\,
m_{r+1+t}\bigl(
\operatorname{id}^{\otimes r} \otimes m_s
\otimes \operatorname{id}^{\otimes t}\bigr)
= 0,
\]
so all contributions vanish. In the chiral setting, the
operadic compositions still carry block-substitution of spectral
variables
(Definition~\ref{def:thqg-chiral-endomorphism-operad}(b)), and
associativity of block-substitution gives the same cancellation
pattern as in the classical brace algebra.

\emph{Step~2: Pre-Lie identity.}
The single-brace product $f \bullet g := f\{g\}$ satisfies the
\emph{right pre-Lie identity}:
$(f \bullet g) \bullet h - f \bullet (g \bullet h)
= (-1)^{\|g\|\,\|h\|}\bigl[(f \bullet h) \bullet g
- f \bullet (h \bullet g)\bigr]$.
This is proved by expanding both sides using the partial
compositions~\eqref{eq:thqg-block-composition}: the left side
involves summing over pairs of insertion positions $(i, j)$ in
$f$, and the right side involves the same sum with $g$ and $h$
interchanged. The difference telescopes because
block-substitution is \emph{sequential}: inserting $g$ then $h$
into non-overlapping slots of $f$ produces the same spectral
variable assignment regardless of order, while inserting into
overlapping slots is governed by nested associativity
(Lemma~\ref{lem:partial-comp-assoc}(b) of
\S\ref{subsec:chiral-endo-operad}).
The pre-Lie identity implies the Jacobi identity for the
Gerstenhaber bracket $[f, g] = f\{g\} - (-1)^{\|f\|\,\|g\|} g\{f\}$.

\emph{Step~3: Higher brace identities and Leibniz rule.}
The multi-fold brace operations satisfy the higher pre-Lie
(brace) identities by the same sequential/nested argument. The
Leibniz rule $\delta(f\{g_1, \ldots, g_r\})
= (\delta f)\{g_1, \ldots, g_r\}
+ \sum_{s} \pm f\{g_1, \ldots, \delta g_s, \ldots, g_r\}
+ \text{boundary terms}$ follows because $\delta = [m, -]$
and $m$ has cochain degree~$1$, so $[m, -]$ acts as a
graded derivation of the brace operations. The boundary terms
are the braces involving $m$ insertions, which cancel pairwise
by the brace form $m\{m\} = 0$ of the $A_\infty$ identities.

\emph{Step~4: Gerstenhaber on cohomology.}
The cup product $f \smile g$ (concatenation composition) and the
Gerstenhaber bracket $[f, g]$ both descend to cohomology because
$\delta$ is a derivation of both. The resulting structure
$(H^*, \smile, [-,-])$ satisfies the Gerstenhaber axioms
(graded commutativity of $\smile$, graded Jacobi for $[-,-]$,
Leibniz compatibility) because these are identities in the brace
dg algebra. The full proof with explicit signs in the chiral
setting is given in
\S\ref{subsec:brace-dg-algebra}
(Theorem~\ref{thm:brace-dg-algebra},
Definition~\ref{def:chiral-hochschild-cochain-brace}).
\end{proof}

\begin{definition}[Chiral derived center]
\label{def:thqg-chiral-derived-center}
\index{derived center!chiral|textbf}
The \emph{cochain-level chiral derived center} of~$\cA$ is the
brace dg algebra
\[
\mathbf{Z}^{\mathrm{der}}_{\mathrm{ch}}(\cA)
:=
\bigl(\mathcal{C}^{\bullet}_{\mathrm{ch}}(\cA,\cA),\delta\bigr).
\]
Its cohomology is denoted
\begin{equation}\label{eq:thqg-derived-center}
\mathcal{Z}^{\mathrm{der}}_{\mathrm{ch}}(\cA)
\;:=\;
H^*\!\bigl(\mathcal{C}^{\bullet}_{\mathrm{ch}}(\cA, \cA),\, \delta\bigr).
\end{equation}
By Theorem~\ref{thm:thqg-brace-dg-algebra}(iv), this is a
Gerstenhaber algebra. The universal closed-sector actor in the
Swiss-cheese theorem is the dg object
$\mathbf{Z}^{\mathrm{der}}_{\mathrm{ch}}(\cA)$; its on-shell
cochain closed fields are the classes in
$\mathcal{Z}^{\mathrm{der}}_{\mathrm{ch}}(\cA)$.
This is a cochain object. The annulus trace is the cyclic Hochschild
chain complex, and topological Hochschild homology is the
spectrum-valued trace of an $E_1$-ring. A physical bulk
factorization algebra of a holomorphic--topological theory compares
to this object after taking its boundary $E_2$ local-operator shadow
and supplying a brace-compatible quasi-isomorphism.
On the Koszul locus at generic level with the
$E_\infty$-chiral PBW completion, Theorem~H
\textup{(}Theorem~\ref{thm:hochschild-polynomial-growth}\textup{)}
gives
$\mathcal{Z}^{\mathrm{der}}_{\mathrm{ch}}(\cA)^n = 0$
for $n \notin \{0, 1, 2\}$, so the derived center is a
three-term Gerstenhaber algebra:
\begin{equation}\label{eq:thqg-derived-center-three-term}
\mathcal{Z}^{\mathrm{der}}_{\mathrm{ch}}(\cA)
\;=\;
\mathcal{Z}^0 \;\oplus\; \mathcal{Z}^1 \;\oplus\; \mathcal{Z}^2,
\end{equation}
with the Lie bracket of degree $-1$ giving
$[\cdot, \cdot] \colon
\mathcal{Z}^1 \otimes \mathcal{Z}^1 \to \mathcal{Z}^1$
(a Lie algebra on bulk derivations) and
$[\cdot, \cdot] \colon
\mathcal{Z}^1 \otimes \mathcal{Z}^2 \to \mathcal{Z}^2$
(the action of derivations on obstructions). No concentration
claim is made here for critical level, minimal/admissible level,
or non-PBW $E_1$-chiral inputs.
\end{definition}

\begin{definition}[Genus-$g$ chiral derived center]
\label{def:genus-g-derived-center}
\index{derived center!genus $g$|textbf}
For a curved $A_\infty$-chiral algebra $(\cA, \{m_k\})$ with
modular characteristic $\kappa(\cA)$
(Definition~\ref{def:modular-characteristic-package}), the
\emph{genus-$g$ chiral derived center} ($g \geq 1$) is
\begin{equation}\label{eq:genus-g-derived-center}
\mathcal{Z}^{\mathrm{der}}_{\mathrm{ch}}(\cA)^{(g)}
\;:=\;
H^*\!\bigl(
\mathcal{C}^{\bullet}_{\mathrm{ch}}(\cA, \cA)
\otimes_{\mathcal{O}_{\overline{\mathcal{M}}_g}}
\Omega^{\bullet}_{\overline{\mathcal{M}}_g}(\log D),
\;\Dg{g}
\bigr),
\end{equation}
where $D \subset \overline{\mathcal{M}}_g$ is the boundary divisor
and the total corrected differential
\begin{equation}\label{eq:Dg-decomposition}
\Dg{g} \;=\; \dfib + \nabla^{\mathrm{GM}}
\end{equation}
has two components: the fiberwise differential $\dfib$ with
$\dfib^{\,2} = \kappa(\cA) \cdot \omega_g \cdot \operatorname{id}$
(Convention~\ref{conv:higher-genus-differentials}), and the
Gauss--Manin connection $\nabla^{\mathrm{GM}}$ on the family of
bar complexes over $\overline{\mathcal{M}}_g$. The total
differential satisfies $\Dg{g}^{\,2} = 0$: the Gauss--Manin
connection absorbs the fiberwise curvature
$\kappa(\cA) \cdot \omega_g$.

At genus~$0$, $\overline{\mathcal{M}}_{0,n}$ carries no Hodge
class ($\omega_0 = 0$), the fiberwise differential is strict
($\dfib^{\,2} = 0$), and
$\mathcal{Z}^{\mathrm{der}}_{\mathrm{ch}}(\cA)^{(0)}$ reduces
to the local chiral derived center of
Definition~\ref{def:thqg-chiral-derived-center}.
\end{definition}

% ===================================================================
\subsection{The universal open/closed pair}
\label{subsec:thqg-universal-oc-pair}
\index{universal open/closed pair|textbf}
% ===================================================================

\begin{definition}[Local chiral open/closed pair]
\label{def:thqg-local-oc-pair}
A \emph{local chiral open/closed pair} with open color~$\cA$
is a triple $(\cB, \cA, \{\mu_{p;q}\})$ in the completed
conformal-weight and pole-order topology where:
\begin{enumerate}[label=\textup{(\alph*)}]
\item $\cB$ is a complete continuous brace dg algebra. Equivalently,
 it is the closed $E_2$ colour in the local chiral Swiss-cheese
 model, with the $E_2$ structure presented by braces.
\item $\cA$ is a complete $A_\infty$-chiral algebra, the open
 $E_1$ colour, whose structure element
 $m=\sum_{n\geq 0}m_n$ is a Maurer--Cartan element in the
 completed chiral endomorphism brace Lie algebra.
\item the mixed operations
 $\mu_{p;q}\colon \cB^{\otimes p}\otimes A^{\otimes q}
 \to A((\lambda_1))\cdots((\lambda_{p+q-1}))$
 are continuous maps of degree $2-p-q$, have finite pole order on
 each completed finite piece, and satisfy the local chiral
 Swiss-cheese codimension-one identities of
 Definition~\ref{def:local-swiss-cheese-pair}.
\end{enumerate}
These identities imply that the one-closed-input operations define
a continuous
brace dg algebra map
\[
\iota_\mu\colon \cB\to
\mathcal{C}^{\bullet}_{\mathrm{ch}}(\cA,\cA),
\qquad
\iota_\mu(b)_n=(-1)^{|b|}\mu_{1;n}(b;-),
\]
encoding the bulk-to-boundary coupling.
Conversely, a continuous brace dg algebra map
$\cB\to\mathcal{C}^{\bullet}_{\mathrm{ch}}(\cA,\cA)$ determines
the local mixed action by iterated brace evaluation.
A \emph{morphism} of local chiral open/closed pairs
$(\cB, \cA, \mu) \to (\cB', \cA, \mu')$ (with fixed open color)
is a brace dg algebra map $\phi \colon \cB \to \cB'$ such that
$\iota_{\mu'} \circ \phi = \iota_\mu$.
\end{definition}

\begin{theorem}[Universal open/closed pair (chiral Swiss-cheese theorem);
\ClaimStatusProvedHere]
\label{thm:thqg-swiss-cheese}
\index{Swiss-cheese theorem!chiral|textbf}
\textit{Type signature.}
Quadrant: Open. Presentation: chain-level brace dg algebra.
Level: $1 \to 3$ (chart $\to$ derived chiral centre).
Hypothesis package: complete curved $A_\infty$-chiral algebra
$(A, \{m_k\})$ with arity and pole-order filtrations complete and
separated; chiral Maurer--Cartan equation $m\{m\} = 0$;
brace dg algebra structure
\textup{(}Theorem~\ref{thm:thqg-brace-dg-algebra}\textup{)}.

\medskip
Let $(A, \{m_k\})$ be a complete curved $A_\infty$-chiral algebra:
the arity and pole-order filtrations are complete and separated, and
$m=\sum_{k\geq 0}m_k$ satisfies the completed chiral
Maurer--Cartan equation $m\{m\}=0$.
The pair
\begin{equation}\label{eq:thqg-universal-oc-pair}
\mathcal{U}(\cA)
\;:=\;
\bigl(\mathcal{C}^{\bullet}_{\mathrm{ch}}(\cA, \cA),\;
\cA,\; \{\mu^{\mathrm{univ}}_{p;q}\}\bigr)
\end{equation}
is the \emph{terminal} \textup{(}final\textup{)} local chiral
Swiss-cheese pair over~$\cA$
\textup{(}Definition~\textup{\ref{def:thqg-local-oc-pair}}\textup{)}:
for every local chiral open/closed pair
$(\cB, \cA, \{\mu_{p;q}\})$ over~$\cA$, there exists a
\emph{unique} morphism
\[
\Phi \colon (\cB, \cA, \{\mu_{p;q}\})
\;\longrightarrow\; \mathcal{U}(\cA)
\]
of local chiral Swiss-cheese pairs.
Here the universal mixed operations
$\mu^{\mathrm{univ}}_{p;q} \colon
(\mathcal{C}^{\bullet}_{\mathrm{ch}})^{\otimes p}
\otimes A^{\otimes q}
\to A((\lambda_1))\cdots((\lambda_{p+q-1}))$
are defined by iterated brace evaluation of chiral Hochschild
cochains on algebra elements. Passing to cohomology gives the
on-shell Gerstenhaber algebra
$\mathcal{Z}^{\mathrm{der}}_{\mathrm{ch}}(\cA)$, but terminality
is a cochain-level statement.
\end{theorem}

\begin{proof}
The morphism $\Phi \colon \cB \to
\mathcal{C}^{\bullet}_{\mathrm{ch}}(\cA, \cA)$ is constructed
by evaluating closed observables on the open sector: for
$b \in \cB$, define
\[
\Phi(b)_n(a_1, \ldots, a_n)
\;:=\;
(-1)^{|b|}\,\mu_{1;n}(b;\, a_1, \ldots, a_n),
\]
where $\mu_{1;n}$ is the mixed operation of the Swiss-cheese pair.
This gives a well-defined element $\Phi(b) \in
\mathcal{C}^{\bullet}_{\mathrm{ch}}(\cA, \cA)$ because
$\mu_{1;n}$ is $\mathbb{C}$-multilinear in the $a_i$ with values
in $A((\lambda_1))\cdots((\lambda_{n-1}))$ by the sesquilinearity
axiom.

The Swiss-cheese codimension-$1$ identities for the pair
$(\cB, \cA, \mu)$ imply:
\begin{enumerate}[label=\textup{(\roman*)}]
\item \emph{Chain map}: $\Phi \circ \delta_\cB = \delta \circ \Phi$
 (from the $(p,q) = (1,n)$ identity);
\item \emph{Brace compatibility}: $\Phi(b\{b_1, \ldots, b_r\})
 = \Phi(b)\{\Phi(b_1), \ldots, \Phi(b_r)\}$
 (from the $(p, q) = (1+r, N)$ identity).
\end{enumerate}
Uniqueness: any Swiss-cheese morphism
$\Psi \colon \cB \to \mathcal{C}^{\bullet}_{\mathrm{ch}}$
must intertwine the mixed operations, forcing
$\Psi(b)_n = (-1)^{|b|}\mu_{1;n}(b; -)$ by evaluation, so
$\Psi = \Phi$.
The full proof with explicit signs and spectral-variable
conventions is given in
Theorem~\ref{thm:chiral-deligne-tamarkin}
(\S\ref{subsec:local-swiss-cheese}).
\end{proof}

\begin{proposition}[Typed physical-bulk comparison;
\ClaimStatusProvedHere]
\label{prop:thqg-universal-action-not-reconstruction}
\index{physical bulk!separation from derived center|textbf}
\index{derived center!universal action versus reconstruction|textbf}
Let $T$ be a holomorphic--topological theory on a half-space with
boundary chiral algebra~$\cA$, and let
$\operatorname{Obs}^{\mathrm{bulk}}(T)$ denote its physical bulk
factorization algebra. A comparison between the physical bulk and
the cochain derived center consists of three typed data:
\begin{enumerate}[label=\textup{(\roman*)}]
\item a boundary identification
 $\operatorname{Obs}^{\mathrm{bdry}}(T)\simeq \cA$ as an
 $A_\infty$-chiral algebra;
\item a boundary local-operator shadow
 $\operatorname{Loc}_{\partial}
 \operatorname{Obs}^{\mathrm{bulk}}(T)$ carrying a continuous
 $E_2$ or brace algebra structure;
\item a continuous brace dg algebra map
 \[
 \beta_T\colon
 \operatorname{Loc}_{\partial}
 \operatorname{Obs}^{\mathrm{bulk}}(T)
 \longrightarrow
 \mathcal{C}^{\bullet}_{\mathrm{ch}}(\cA,\cA)
 \]
 whose induced mixed operations on~$\cA$ agree with the
 bulk-to-boundary operations of~$T$.
\end{enumerate}
The chiral Swiss-cheese theorem supplies the unique map
$\beta_T$ once the physical bulk acts on the boundary. A
quasi-isomorphism condition on~$\beta_T$ is the extra hypothesis
which identifies the boundary $E_2$ shadow of the physical bulk with
$\mathbf{Z}^{\mathrm{der}}_{\mathrm{ch}}(\cA)$. The theorem itself
classifies universal local actions on~$\cA$.
\end{proposition}

\begin{proof}
The boundary action of
$\operatorname{Loc}_{\partial}\operatorname{Obs}^{\mathrm{bulk}}(T)$
on~$\cA$ is precisely a local chiral Swiss-cheese pair. Applying
Theorem~\ref{thm:thqg-swiss-cheese} gives a unique brace dg algebra
map to
$\mathcal{C}^{\bullet}_{\mathrm{ch}}(\cA,\cA)$. This uniqueness is
internal to the category of local actions. A physical bulk factorization
algebra carries descent, locality, and topological-direction data in
addition to its boundary local-operator shadow. Recovering those data
requires the specified comparison map and a proof that it is a
quasi-isomorphism in the ambient completion.
\end{proof}

\begin{remark}[Bulk, Verdier dual, and bar-cobar inverse]
\label{rem:thqg-oc-physical}
In the holomorphic-topological framework of
Costello--Dimofte--Gaiotto, the boundary of a 3d HT theory is an
$A_\infty$-chiral algebra $\cA$, and the bulk local operators map to
the derived center of $\cA$. Theorem~\ref{thm:thqg-swiss-cheese}
identifies the universal acting cochain object with
$\mathbf{Z}^{\mathrm{der}}_{\mathrm{ch}}(\cA)$.
This is independent of the two bar constructions in the spine of
the monograph:
\[
\Omega^{\mathrm{ch}}\barB^{\mathrm{ch}}(\cA)\simeq \cA
\quad\text{is bar-cobar inversion,}
\]
\[
\mathbb{D}_{\Ran}\barB_X(\cA)\simeq \cA^!_\infty
\quad\text{is Verdier duality,}
\]
whereas
$\mathbf{Z}^{\mathrm{der}}_{\mathrm{ch}}(\cA)$ is governed by
chiral Hochschild cochains. Presenting the Koszul dual line algebra
as a Hochschild center requires a compact-generator Hochschild
comparison for the relevant category of boundary conditions.
\end{remark}

\begin{proposition}[Mixed sector encodes bulk-to-boundary module structure;
\ClaimStatusProvedHere]
\label{prop:mixed-sector-bulk-boundary}
\index{Swiss-cheese operad!mixed sector|textbf}
\index{bulk-to-boundary!mixed sector identification}
The mixed sector of the Koszul dual cooperad
$\mathsf{SC}^{\mathrm{ch,top},!}$ at bi-degree $(k,m)$,
$k \geq 1$, with
$\dim = (k{-}1)!\,\binom{k+m}{m}$
\textup{(}Proposition~\textup{\ref{prop:sc-koszul-dual-three-sectors}}\textup{)},
encodes the bulk-to-boundary module structure: the iterated
brace evaluation
$\mu_{k;m}\colon
\bigl(\mathbf{Z}^{\mathrm{der}}_{\mathrm{ch}}(\cA)\bigr)^{\otimes k}
\otimes \cA^{\otimes m} \to \cA$
by which the chiral derived center acts on the boundary algebra.
At $(k{=}1,\,m{=}n)$, this recovers the $n$-fold brace
evaluation with $\dim = n{+}1$ counting the insertion positions
for one bulk operator among $n$ ordered boundary points.
\end{proposition}

\begin{proof}
The mixed operations
$\mathsf{SC}^{\mathrm{ch,top}}(\mathsf{c}^k,\mathsf{o}^m;\mathsf{o})
= \FM_k(\mathbb{C}) \times E_1(m)$
parametrize $k$ closed (bulk) inputs and $m$ open (boundary)
inputs producing one open output. By
Theorem~\ref{thm:thqg-swiss-cheese}, the terminal open/closed
pair
$(\mathbf{Z}^{\mathrm{der}}_{\mathrm{ch}}(\cA),\,\cA)$
realizes these operations as iterated brace evaluations. The
K\"unneth factorization
$\operatorname{Lie}^c(k) \otimes \operatorname{Sh}(k,m)$
(where $\operatorname{Sh}(k,m) = \binom{k+m}{m}$ counts
$(k,m)$-shuffles) gives the dimension.
\end{proof}

\begin{proposition}[Yangian coproducts require line-category data;
\ClaimStatusProvedHere]
\label{prop:thqg-swiss-cheese-no-yangian-coproduct}
\index{Yangian!requires line-category data|textbf}
\index{coproduct!requires line-category data|textbf}
An algebra over $\mathsf{SC}^{\mathrm{ch,top}}$ supplies operations
of types
\[
(\mathsf{c}^k\to\mathsf{c}),\qquad
(\mathsf{o}^m\to\mathsf{o}),\qquad
(\mathsf{c}^k,\mathsf{o}^m\to\mathsf{o}).
\]
A coproduct on~$\cA^!$, a bialgebra structure on the line algebra,
or a Yangian/RTT presentation requires additional data: a monoidal
tensor product on the line category $\mathsf{Line}(\cA)$, a compact
generator identifying its endomorphism algebra with the
Verdier/Koszul dual line algebra~$\cA^!$, a coproduct induced by
tensoring line operators, and an RTT or KZ flatness comparison
identifying the binary collision residue $r(z)$ with the relevant
Yangian $R$-matrix.
\end{proposition}

\begin{proof}
The Swiss-cheese operad and its Koszul dual cooperad organize
many-input operations and the corresponding bar co-operations in the
operadic resolution. A Yangian coproduct is a one-input,
many-output structure on a line-operator algebra. Its construction
uses the monoidal line category rather than the mixed operations
$\mu_{k;m}$ or the dimension formula in
Proposition~\ref{prop:sc-koszul-dual-three-sectors}. The bar
construction produces the coalgebra $B^{\mathrm{ch}}(\cA)$, and
Verdier duality produces the algebra~$\cA^!$ on the finite-type or
completed continuous Koszul locus. A bialgebra or Yangian structure
on~$\cA^!$ is supplied by the monoidal line category and the
separate RTT/KZ comparison data named in the statement.
\end{proof}

\begin{remark}[Mixed sector and multi-weight channels]
\label{rem:mixed-sector-genus-independent}
The product decomposition
$C_*(\FM_{k|m}(\Sigma_g, \partial))
\simeq C_*(\FM_k(\Sigma_g)) \times C_*(E_1(m))$
(Volume~II, Proposition~\ref{V2-prop:mixed-product-decomposition} in the
modular Swiss-cheese chapter) holds at all genera.
The open colour $E_1(m)$ is genus-independent: the interval~$\mathbb{R}$
has no genus. All genus dependence lives in the closed factor. The
curvature $\kappa(\cA) \cdot \omega_g$ is closed-sector data, and
genus-raising operations act on closed inputs. The mixed channels in
the multi-weight correction
$\delta F_g^{\mathrm{cross}}$ are mixed-\emph{weight} channels
\textup{(}closed-sector endomorphism decorations\textup{)} rather
than mixed-\emph{sector} bulk-to-boundary $\mathsf{SC}$ operations.
\end{remark}

% ===================================================================
\subsection{Restriction from global to local: the bordered FM bridge}
\label{subsec:thqg-bordered-fm-bridge}
\index{bordered FM bridge|textbf}
% ===================================================================

The bordered FM compactification
$\overline{C}{}^{\,\mathrm{oc}}_{k,\vec{m}}(X, D, \tau)$
of Construction~\ref{constr:bordered-fm} provides the global
geometric stage. Its formal-disk restriction has two outputs:
the closed local chart is the chiral Hochschild cochain complex
of~$\cA$, while the open categorical chart is the Morita
Hochschild center of the boundary category.

\begin{theorem}[Local-global bridge; \ClaimStatusProvedHere]
\label{thm:thqg-local-global-bridge}
\index{local-global bridge!open/closed|textbf}
Let $(X, D, \tau)$ be a tangential log curve with a single
puncture $D = \{p\}$, and let $\cA$ be a chiral algebra on~$X$
with module $\mathcal{M}$ at~$p$. Assume $\mathcal{M}$ admits
a compact generator~$b$ and is Morita equivalent to
$\operatorname{Perf}(A_b)$ for
$A_b:=\operatorname{RHom}_{\mathcal M}(b,b)$. Then:
\begin{enumerate}[label=\textup{(\roman*)}]
\item \emph{Algebraic chart.}
 The boundary algebra
 $A_b := \operatorname{RHom}_{\mathcal{M}}(b, b)$
 is an $A_\infty$-algebra, with operations $\mu_n^{A_b}$ encoding
 the ordered boundary collisions of
 $\operatorname{Conf}^{\mathrm{ord}}_n(I_p)$
 (Definition~\ref{def:mixed-config-space}).

\item \emph{Closed formal chart.}
 The genus-$0$, pure-interior component of the open-closed
 convolution algebra
 $\mathfrak{g}^{\mathrm{oc}}_{\cA, \mathcal{M}}$
 (Definition~\ref{def:oc-modular-convolution}), restricted to
 the formal disk at~$p$, is canonically quasi-isomorphic to the
 algebraic chiral Hochschild cochain complex:
 \begin{equation}\label{eq:thqg-local-global-identification}
 \mathfrak{g}^{\mathrm{oc}}_{\cA, \mathcal{M}}
 \bigl|_{g=0,\, \text{int-only},\, \widehat{D}_p}
 \;\simeq\;
 \mathcal{C}^{\bullet}_{\mathrm{ch}}(\cA, \cA).
 \end{equation}
 This chart sees the closed algebra~$\cA$ and is independent of the
 compact generator~$b$.

\item \emph{Open categorical center.}
 The derived center of the boundary category is computed in the
 compact-generator chart by
 \[
 Z_{\mathrm{der}}(\mathcal M)
 \;\simeq\;
 \operatorname{RHom}_{A_b^e}(A_b,A_b)
 =
 \mathrm{HH}^{\bullet}_{\mathrm{alg}}(A_b).
 \]
 Here $A_b^e=A_b\otimes A_b^{\mathrm{op}}$.
 This quasi-isomorphism type depends only on the Morita
 equivalence class of~$\mathcal M$ and is independent of the choice
 of~$b$.
 If the boundary endomorphism algebra is the formal chiral
 boundary algebra induced from~$\cA$, the comparison of
 Proposition~\ref{prop:geometric-algebraic-hochschild} identifies
 this algebraic center with the corresponding chiral Hochschild
 cochain complex.
\end{enumerate}
\end{theorem}

\begin{proof}
(i) The ordered boundary configuration space
$\operatorname{Conf}^{\mathrm{ord}}_n(I_p)$ is the interior of
the simplex $\Delta^{n-1}$
(Remark~\ref{rem:boundary-ordering-associahedron}).
Its compactification is the associahedron $K_n$, whose
codimension-one faces parametrize consecutive-block compositions.
The boundary corrections give the $A_\infty$ operations $\mu_n^{A_b}$
by the standard Stasheff construction on
$C_*(\overline{K}_n) \otimes \operatorname{End}(b)$.

(ii) The interior configuration space $C_k(X \setminus D)$,
restricted to the formal disk $\widehat{D}_p$, is the
configuration space of $k$ points in a formal punctured disk.
After choosing a local coordinate, this identifies with the
space of $k$ spectral variables $\lambda_1, \ldots, \lambda_{k-1}$
(the relative positions). The FM compactification on this formal
disk gives exactly the operadic compositions of
$\mathcal{E}\!nd^{\mathrm{ch}}_A$
(Definition~\ref{def:thqg-chiral-endomorphism-operad}):
codimension-one faces correspond to operadic insertions with
block-substitution. The identification
$C_*(\overline{C}_k(\widehat{D}_p))
\otimes \operatorname{Hom}(A^{\otimes k}, A)
\cong \mathcal{E}\!nd^{\mathrm{ch}}_A(k)$
at each degree gives~\eqref{eq:thqg-local-global-identification}.

(iii) Standard Morita theory: if $b' \in \mathcal{M}$ is
another compact generator with
$\mathcal{M} \simeq \operatorname{Perf}(A_b)
\simeq \operatorname{Perf}(A_{b'})$,
then Theorem~\ref{thm:derived-center-hochschild} gives
$Z_{\mathrm{der}}(\mathcal M)\simeq \operatorname{RHom}_{A_b^e}(A_b,A_b)$
and the analogous expression for~$A_{b'}$. The bimodule
implementing the Morita equivalence carries one Hochschild
cochain complex to the other.
\end{proof}

\begin{corollary}[The intrinsic categorical center]
\label{cor:thqg-intrinsic-bulk}
The \emph{intrinsic categorical center} of the open sector is the
Morita-invariant object
\begin{equation}\label{eq:thqg-intrinsic-bulk}
Z_{\mathrm{der}}(\mathcal{M})
\;:=\;
\operatorname{RHom}_{\operatorname{Fun}(\mathcal{M},
\mathcal{M})}(\mathrm{Id},\, \mathrm{Id}),
\end{equation}
which in a chart $\mathcal{M} \simeq \operatorname{Perf}(A_b)$
is computed by $\mathrm{HH}^{\bullet}_{\mathrm{alg}}(A_b)$.
This is the center of the boundary category. It becomes the closed
chiral bulk
$\mathbf{Z}^{\mathrm{der}}_{\mathrm{ch}}(\cA)$ only under the
completed formal chiral boundary comparison: $A_b$ must be the
completed formal chiral boundary algebra induced from~$\cA$, and
the logarithmic finite-piece hypotheses of
Proposition~\ref{prop:geometric-algebraic-hochschild} must hold.
Under these hypotheses the geometric--algebraic comparison identifies
$\mathrm{HH}^{\bullet}_{\mathrm{alg}}(A_b)$ with the corresponding
chiral cochain model
$\mathcal{C}^{\bullet}_{\mathrm{ch}}(A_b, A_b)$. The concentration
in degrees~$\{0,1,2\}$ transfers only on the Koszul/PBW locus of
Theorem~H, with the same generic-level exclusions as in
Definition~\ref{def:thqg-chiral-derived-center}.
\end{corollary}

\begin{theorem}[Separation of Hochschild, trace, THH, BV, and
Koszul duality; \ClaimStatusProvedHere]
\label{thm:thqg-hochschild-register-separation}
\index{Hochschild theories!open/closed separation|textbf}
\index{trace!open/closed separation|textbf}
\index{topological Hochschild homology!separation from chiral center|textbf}
Let $\cA$ be a chiral algebra on a smooth curve, let
$\mathcal{M}$ be a boundary dg category with compact generator~$b$,
and set $A_b=\operatorname{RHom}_{\mathcal{M}}(b,b)$. The local
open/closed construction contains four different Hochschild-type
objects and two comparison mechanisms:
\begin{enumerate}[label=\textup{(\roman*)}]
\item \emph{Chiral cochains.}
 The chiral derived center is
 \[
 \mathbf{Z}^{\mathrm{der}}_{\mathrm{ch}}(\cA)
 =
 \mathcal{C}^{\bullet}_{\mathrm{ch}}(\cA,\cA),
 \qquad
 \delta=[m,-].
 \]
 It is the terminal closed colour acting on the open chiral algebra
 in Theorem~\ref{thm:thqg-swiss-cheese}.

\item \emph{Categorical cochains.}
 The derived center of the boundary category is
 \[
 Z_{\mathrm{der}}(\mathcal{M})
 =
 \operatorname{RHom}_{\operatorname{Fun}(\mathcal M,\mathcal M)}
 (\mathrm{Id},\mathrm{Id})
 \simeq
 \operatorname{RHom}_{A_b^e}(A_b,A_b),
 \]
 by Theorem~\ref{thm:derived-center-hochschild}. This is a Morita
 invariant of $\mathcal M$. It becomes a chiral invariant of~$\cA$
 only after the compact-generator chart is identified with the
 formal chiral boundary algebra by
 Proposition~\ref{prop:geometric-algebraic-hochschild}.

\item \emph{Categorical chains.}
 The annulus trace is
 \[
 \operatorname{Tr}(\mathcal M)
 =
 \int_{S^1_p}\mathcal M
 \simeq
 \mathrm{HH}_{\bullet,\mathrm{cat}}(\mathcal M)
 \simeq
 A_b\otimes^{\mathbf L}_{A_b^e}A_b .
 \]
 It is the categorical chain trace. A comparison to
 $\mathbf{Z}^{\mathrm{der}}_{\mathrm{ch}}(\cA)$ requires the
 boundary local-operator map.

\item \emph{Topological Hochschild homology.}
 If $A_b$ is spectrified to an $E_1$-algebra over the
 Eilenberg--Mac Lane spectrum $Hk$, written $H A_b$, then
 \[
 \operatorname{THH}_{Hk}(H A_b)
 =
 H A_b
 \otimes^{\mathbf L}_{H A_b\wedge_{Hk}H A_b^{\mathrm{op}}}
 H A_b
 \]
 is an $S^1$-equivariant spectrum. It is the spectral lift of the
 circle trace. The chiral Hochschild cochain complex is the separate
 cochain slot. A rational or Eilenberg--Mac Lane comparison, when
 supplied, recovers the algebraic cyclic bar complex after applying
 chains; such a comparison is additional to
 Theorem~\ref{thm:thqg-swiss-cheese}.

\item \emph{Calabi--Yau duality and BV.}
 If $\mathcal{M}$ is a smooth proper $d$-Calabi--Yau category with
 cyclic class $\eta_{\mathcal M}$, then cap product gives a
 Poincar\'e duality map
 \[
 \mathrm{PD}_{\eta_{\mathcal M}}\colon
 \mathrm{HH}_{\bullet,\mathrm{cat}}(\mathcal M)
 \xrightarrow{\;\sim\;}
 \mathrm{HH}^{d-\bullet}_{\mathrm{cat}}(\mathcal M).
 \]
 Under this extra datum the Connes operator on chains transports
 to a BV operator on cochains. The Calabi--Yau class
 $\eta_{\mathcal M}$ is the chain-to-cochain comparison datum between
 the annulus trace and the derived center.

\item \emph{Koszul dual algebra.}
The homotopy Koszul dual algebra
$\cA^!_\infty\simeq \mathbb{D}_{\Ran}\barB_X(\cA)$ is produced
from the bar coalgebra by Verdier duality. On the finite-type
Koszul locus its cohomology gives the strict algebra
$\cA^!=(\cA^i)^\vee$. This is the line-operator algebra in the
modular Koszul datum. The derived-center and categorical-trace
slots are respectively
$\mathbf{Z}^{\mathrm{der}}_{\mathrm{ch}}(\cA)$ and
$\operatorname{Tr}(\mathcal M)$.
\end{enumerate}
A derivation of the closed-sector scalar formula
$F_g=\kappa(\cA)\lambda_g$ on the uniform-weight lane from the
categorical trace theorem~\eqref{eq:thqg-annulus-trace} or from the
spectral object $\operatorname{THH}$ requires an explicit
open/closed MC element, a mixed bulk-to-boundary map, and, when
chains are converted to cochains, a specified Calabi--Yau duality
class.
\end{theorem}

\begin{proof}
Items (i)--(iii) are the definitions and identifications in
Definition~\ref{def:thqg-chiral-derived-center},
Theorem~\ref{thm:derived-center-hochschild}, and
Theorem~\ref{thm:thqg-annulus-trace}. Item~(iv) is the definition
of topological Hochschild homology for an $E_1$-ring spectrum, with
the input and output type fixed by
Convention~\ref{conv:three-hochschild}. Item~(v) is the standard
Calabi--Yau Hochschild Poincar\'e duality: the cyclic class
identifies the diagonal bimodule with its shifted dual, carrying
chains to cochains and transporting the Connes operator to BV.
Item~(vi) is the Verdier construction of the Koszul dual algebra
already separated in Remark~\ref{rem:thqg-oc-physical}.

The final assertion follows from type. The annulus trace lives in
categorical Hochschild homology of $\mathcal M$; the closed formula
$F_g=\kappa(\cA)\lambda_g$ is the scalar reduction of the
closed-sector MC element in the modular convolution algebra. These
objects have different inputs and different targets. A comparison is
extra structure, supplied only by the mixed sector of
$\Theta^{\mathrm{oc}}_\cA$ together with the duality data named
above.
\end{proof}

% ===================================================================
\subsection{The annulus trace: first modular shadow of the open sector}
\label{subsec:thqg-annulus-trace}
\index{annulus trace|textbf}
\index{categorical Hochschild homology!as open trace|textbf}
% ===================================================================

\begin{theorem}[Annulus trace theorem; \ClaimStatusProvedHere]
\label{thm:thqg-annulus-trace}
Let $(X, D, \tau)$ be a tangential log curve with a single
puncture~$p$, and let $\mathcal{M}$ be the boundary category
\textup{(}an $E_1$-category, i.e.\ a monoidal dg-category whose
monoidal structure is locally constant along the boundary interval
$I_p$\textup{)}. There is a canonical quasi-isomorphism
\begin{equation}\label{eq:thqg-annulus-trace}
\int_{S^1_p} \mathcal{M}
\;\simeq\;
\mathrm{HH}_{\bullet,\mathrm{cat}}(\mathcal{M}).
\end{equation}
In a chart $\mathcal{M} \simeq \operatorname{Perf}(A_b)$
\textup{(}Theorem~\textup{\ref{thm:thqg-local-global-bridge}(i)}\textup{)},
the right side is computed by the classical algebraic Hochschild
homology of the $A_\infty$-algebra $A_b$:
\begin{equation}\label{eq:thqg-hochschild-bar-model}
\mathrm{HH}_{\bullet,\mathrm{alg}}(A_b)
\;=\; H_*\!\bigl(B^{\mathrm{cyc}}(A_b), \, d_{\mathrm{cyc}}\bigr),
\end{equation}
where $B^{\mathrm{cyc}}(A_b) = \bigoplus_{n \ge 0}
(s A_b)^{\otimes n} \otimes_{(\mathbb{Z}/n\mathbb{Z})} A_b$
is the cyclic bar complex.
\end{theorem}

\begin{proof}
The argument applies the Ayala--Francis excision framework
for factorization homology~\cite{AF15} to the boundary circle
$S^1_p$ viewed as a $1$-manifold on which $\mathcal{M}$ is a
locally constant factorization category (by axiom~(iv) of
Definition~\ref*{V2-def:oc-factorization-category}). The same
excision argument appears for locally constant factorization
algebras in
Theorem~\ref{thm:circle-fh-hochschild}
(\S\ref{sec:annulus-hochschild-chains}); the present proof
adapts it to the boundary-circle setting of the open/closed
factorization category, retaining full detail for self-containment.

The proof has five parts: an excision decomposition of
$S^1_p$, identification of the resulting bar-type complex,
verification of the differential, the passage to cyclic
invariants, and Morita functoriality.

\medskip
\noindent\textit{Step~1: Excision decomposition of $S^1_p$.}
Choose two distinct points $\star_+, \star_- \in S^1_p$ with
$\star_+ \ne \star_-$, and let
\begin{equation}\label{eq:thqg-excision-cover}
U_+ := S^1_p \setminus \{\star_-\},
\qquad
U_- := S^1_p \setminus \{\star_+\}.
\end{equation}
Each $U_\pm$ is diffeomorphic to $\mathbb{R}$ and the intersection
$U_+ \cap U_- = S^1_p \setminus \{\star_+, \star_-\}$ has two
connected components, each diffeomorphic to an open interval.
The pair $\{U_+, U_-\}$ forms a Weiss cover of $S^1_p$ in the
sense of~\cite{AF15}: every finite subset of $S^1_p$ is contained
in some $U_\alpha$. The excision theorem for factorization
homology \cite[Theorem~3.24]{AF15} gives a canonical equivalence
\begin{equation}\label{eq:thqg-excision-formula}
\int_{S^1_p} \mathcal{M}
\;\simeq\;
\int_{U_+} \mathcal{M}
\;\underset{\int_{U_+ \cap U_-} \mathcal{M}}
{\overset{\mathbf{L}}{\otimes}}\;
\int_{U_-} \mathcal{M}\,.
\end{equation}
This is a derived tensor product over the factorization homology
of the intersection, with the two maps given by the restriction
functors $U_\pm \hookleftarrow U_+ \cap U_-$.

\medskip
\noindent\textit{Step~2: Evaluation on intervals.}
Since $\mathcal{M}$ is an $E_1$-category (locally constant
factorization category on the $1$-manifold $S^1_p$, axiom~(iv)
of Definition~\ref*{V2-def:oc-factorization-category}),
factorization homology over any interval
$J \cong \mathbb{R}$ recovers the underlying category:
\begin{equation}\label{eq:thqg-interval-fh}
\int_J \mathcal{M} \;\simeq\; \mathcal{M}.
\end{equation}
This is because $\mathbb{R}$ is contractible, so the
factorization homology collapses: the constant cosheaf on
$\operatorname{Ran}(\mathbb{R})$ is computed by global
sections~\cite[Proposition~5.4.5.15]{HA},
\cite[Theorem~3.19]{AF15}.

In a chart $\mathcal{M} \simeq \operatorname{Perf}(A_b)$
with compact generator~$b$, the evaluation~\eqref{eq:thqg-interval-fh}
is witnessed by the $A_\infty$-algebra
$A_b = \operatorname{RHom}_{\mathcal{M}}(b, b)$: one has
$\int_\mathbb{R} A_b \simeq A_b$ as an $A_b$-bimodule (the
``diagonal'' bimodule).

\medskip
\noindent\textit{Step~3: Identification of the derived tensor product.}
Applying Step~2 to the three terms
in~\eqref{eq:thqg-excision-formula}:
\begin{align}
\int_{U_+} \mathcal{M} &\simeq A_b
 \quad\text{(right $A_b^e$-module via restriction to }
 U_+ \cap U_-\text{)},
 \label{eq:thqg-U-plus} \\
\int_{U_-} \mathcal{M} &\simeq A_b
 \quad\text{(left $A_b^e$-module)},
 \label{eq:thqg-U-minus} \\
\int_{U_+ \cap U_-} \mathcal{M} &\simeq A_b \otimes A_b^{\mathrm{op}}
 =: A_b^e
 \quad\text{(the enveloping algebra)}.
 \label{eq:thqg-intersection}
\end{align}
For~\eqref{eq:thqg-intersection}: the intersection
$U_+ \cap U_-$ has two connected components, say $V_L$ and $V_R$,
each diffeomorphic to $\mathbb{R}$. Factorization homology is
symmetric monoidal with respect to disjoint union
\cite[Theorem~3.19]{AF15}, so
$\int_{V_L \sqcup V_R} \mathcal{M}
\simeq
\bigl(\int_{V_L} \mathcal{M}\bigr) \otimes
\bigl(\int_{V_R} \mathcal{M}\bigr)
\simeq A_b \otimes A_b$.
The two inclusions $V_L, V_R \hookrightarrow U_+$
(resp.\ $U_-$) compose to give $A_b$ the structure of an
$(A_b, A_b)$-bimodule, equivalently a right (resp.\ left)
$A_b^e$-module via $a \cdot (x \otimes y^{\mathrm{op}}) = x \cdot a
\cdot y$ (resp.\ $(x \otimes y^{\mathrm{op}}) \cdot a = y \cdot a
\cdot x$ with $A_\infty$ corrections). The bimodule structure
on the right side in~\eqref{eq:thqg-U-plus} has the left action
coming from the inclusion $V_L \hookrightarrow U_+$ and the right
action from $V_R \hookrightarrow U_+$; concretely, both act by
the $A_\infty$-multiplication of~$A_b$.

Substituting into~\eqref{eq:thqg-excision-formula}:
\begin{equation}\label{eq:thqg-derived-tensor}
\int_{S^1_p} \mathcal{M}
\;\simeq\;
A_b \;\underset{A_b^e}{\overset{\mathbf{L}}{\otimes}}\; A_b\,.
\end{equation}
This is a derived tensor product of bimodules: $A_b$ is a right
$A_b^e$-module (from the right interval decomposition) and
$A_b$ is a left $A_b^e$-module (from the left interval
decomposition).

\medskip
\noindent\textit{Step~4: The two-sided bar complex resolves the
derived tensor product.}
The derived tensor product~\eqref{eq:thqg-derived-tensor} is computed
by the \emph{two-sided bar complex}
$B(A_b, A_b^e, A_b)$:
\begin{equation}\label{eq:thqg-two-sided-bar}
B(A_b, A_b^e, A_b)_n
\;=\;
A_b \otimes (sA_b)^{\otimes n} \otimes A_b, \qquad n \ge 0.
\end{equation}
The differential $d_{\mathrm{bar}}$ on $B(A_b, A_b^e, A_b)$ is
the sum of the internal $A_\infty$-differential $b'$ and the
simplicial face maps $b''$: for a typical element
$a_0 [a_1 | \cdots | a_n] a_{n+1}$,
\begin{align}
b''(a_0 [a_1 | \cdots | a_n] a_{n+1})
&= \sum_{i=0}^{n-1} (-1)^{\epsilon_i}\,
 a_0 [a_1 | \cdots | m_2(a_i, a_{i+1}) | \cdots | a_n] a_{n+1}
 \notag \\
&\quad + (-1)^{\epsilon_n} m_2(a_0, a_1) [a_2 | \cdots | a_n] a_{n+1}
 \notag \\
&\quad + (-1)^{\epsilon_{n+1}} a_0 [a_1 | \cdots | a_{n-1}]
 m_2(a_n, a_{n+1}),
 \label{eq:thqg-bar-differential}
\end{align}
where $\epsilon_i$ are the standard Koszul signs for desuspended
elements (Appendix~\ref{app:signs}). The higher
$A_\infty$-operations $m_k$ for $k \ge 3$ contribute the
$b'$ terms (the internal differential of the bimodule resolution).

The two-sided bar complex~\eqref{eq:thqg-two-sided-bar} is
identified with the \emph{cyclic bar complex}
$B^{\mathrm{cyc}}(A_b)$: an element
$a_0 [a_1 | \cdots | a_n] a_{n+1}$ in~\eqref{eq:thqg-two-sided-bar}
with $a_0$ and $a_{n+1}$ not suspended is the same datum as
a cyclic word $a_0 \otimes a_1 \otimes \cdots \otimes a_n
\otimes a_{n+1}$ with the identification
$(a_0, a_1, \ldots, a_n, a_{n+1})
\sim
(a_{n+1}, a_0, a_1, \ldots, a_n)$
(cyclic permutation of the unsuspended slots). The
differential~\eqref{eq:thqg-bar-differential} becomes the
standard algebraic Hochschild boundary operator under this identification,
because the last face map wraps $a_n$ past $a_{n+1}$ to compose
with $a_0$; this is the cyclic trace term that
distinguishes the algebraic Hochschild differential from the bar
differential.

Therefore:
\begin{equation}\label{eq:thqg-bar-hochschild-id}
A_b \;\underset{A_b^e}{\overset{\mathbf{L}}{\otimes}}\; A_b
\;\simeq\;
H_*\!\bigl(B^{\mathrm{cyc}}(A_b),\, d_{\mathrm{cyc}}\bigr)
\;=\; \mathrm{HH}_{\bullet,\mathrm{alg}}(A_b).
\end{equation}
The first equivalence is that the two-sided bar complex computes
the derived tensor product (it is a free $A_b^e$-resolution of
$A_b$; see~\cite[\S5.3]{Loday98},
\cite[\S2.2]{LV12}). The second is the
definition of algebraic Hochschild homology in chain complexes.

\medskip
\noindent\textit{Step~5: Morita invariance.}
The quasi-isomorphism type of $\int_{S^1_p} \mathcal{M}$
depends only on the Morita equivalence class
of~$\mathcal{M}$ and is independent of the choice of compact
generator~$b$.
This holds because factorization homology is a functor of the
$E_1$-category~\cite[Theorem~3.19]{AF15}, and Morita equivalent
$E_1$-categories (those with equivalent module categories)
produce canonically quasi-isomorphic factorization homology:
if $b' \in \mathcal{M}$ is another compact generator with
$A_{b'} = \operatorname{RHom}(b', b')$, the Morita bimodule
$M = \operatorname{RHom}(b, b')$ induces a quasi-isomorphism
$\mathrm{HH}_{\bullet,\mathrm{alg}}(A_b)
\simeq \mathrm{HH}_{\bullet,\mathrm{alg}}(A_{b'})$ via the derived trace
$\operatorname{Tr}_M\colon
A_b \otimes_{A_b^e}^{\mathbf{L}} A_b
\xrightarrow{\sim}
A_{b'} \otimes_{A_{b'}^e}^{\mathbf{L}} A_{b'}$
(\cite[Theorem~1.2]{Keller06}).

Combining Steps~1--5:
$\int_{S^1_p} \mathcal{M}
\simeq A_b \otimes_{A_b^e}^{\mathbf{L}} A_b
\simeq \mathrm{HH}_{\bullet,\mathrm{alg}}(A_b)
= \mathrm{HH}_{\bullet,\mathrm{cat}}(\mathcal{M})$.
\end{proof}

\begin{remark}[Relation to the Ayala--Francis excision axiom]
\label{rem:thqg-annulus-excision-principle}
\index{excision!for factorization homology}
The excision decomposition in Step~1 is the instance of the
Ayala--Francis~\cite{AF15} ``$\otimes$-excision'' axiom for
factorization homology applied to the standard Weiss cover of
$S^1$. The key geometric input is that $S^1$ admits a cover by
two contractible open sets whose intersection is the disjoint
union of two contractible open sets, a feature specific to the
circle among closed $1$-manifolds (the unique non-simply-connected
case). This is why factorization homology of an $E_1$-algebra
over $S^1$ produces a two-sided bar construction (two external
modules tensored over an enveloping algebra from two intersection
components), hence the algebraic/categorical Hochschild trace.
After Eilenberg--Mac Lane spectrification the analogous spectral
circle object is $\mathrm{THH}$ of the spectrified $E_1$-algebra.
The present theorem remains in chain complexes.

For higher-dimensional manifolds $M$, the excision axiom still
holds but the resulting complex involves higher-degree factorization
data from the topology of $M$. The present theorem uses only the
$1$-dimensional case, where the open sector is an $E_1$-category
and the computation is purely algebraic.
\end{remark}

\begin{remark}[The annulus as first modular shadow]
\label{rem:thqg-annulus-first-shadow}
The annulus $S^1_p \times [0, 1]$ is the simplest bordered surface
with two boundary components. Its partition function in the
open-closed theory is the \emph{trace} of the open sector:
$Z_{\mathrm{ann}} = \mathrm{Tr}_{\mathcal{M}}(\mathrm{Id})$.
By the annulus trace theorem, this is the categorical Hochschild
trace class corresponding in a chart to the algebraic Hochschild
class $[1] \in \mathrm{HH}_{0,\mathrm{alg}}(A_b)$.
This is the \emph{first modular shadow of the open sector}: a trace
of the open world that precedes the genus-$1$ curvature
$\kappa(\cA) \lambda_1$ in the modular hierarchy. The
Calabi--Yau/cyclic structure on $\mathcal{M}$, when a
$d$-Calabi--Yau class $\eta_{\mathcal M}$ is specified, identifies
$\mathrm{HH}_{\bullet,\mathrm{alg}}(A_b)
\simeq \mathrm{HH}^{d - *}_{\mathrm{alg}}(A_b)$ by Poincar\'e
duality. This is the only canonical chain-to-cochain bridge from
the annulus trace to the derived center.
\end{remark}

\begin{proposition}[Annulus degeneration and the genus-$1$ curvature;
\ClaimStatusProvedHere]
\label{prop:thqg-annulus-degeneration-kappa}
\index{annulus degeneration}
\index{kappa@$\kappa$!from annulus degeneration}
Let\/ $\cA$ be a modular Koszul chiral algebra with boundary
category~$\mathcal{M}$ and annulus trace class
$\operatorname{Tr}_\cA
= [1] \in \mathrm{HH}_{0,\mathrm{cat}}(\mathcal{M})$.
Assume:
\begin{enumerate}[label=\textup{(\roman*)}]
\item $\mathcal M$ is equipped with a cyclic
$d$-Calabi--Yau class $\eta_{\mathcal M}$, so that the two
annulus boundary components can be contracted by the pairing
of Theorem~\ref{thm:thqg-hochschild-register-separation}(v);
\item the open/closed MC datum contains a chain-level
non-separating sewing map
\[
\Delta_{\mathrm{ns}}^{\mathrm{oc}}
\colon
\mathrm{HH}_{\bullet,\mathrm{cat}}(\mathcal M)
\longrightarrow
H^*(\overline{\mathcal{M}}_{1,1})
\]
whose closed projection is the genus-$1$ component
$\Theta_\cA^{(1,0)}$ of the bar-intrinsic MC element;
\item the genus-$1$ closed sector satisfies the scalar
normalization of Theorem~D.
\end{enumerate}
Then the non-separating nodal degeneration of the annulus
$S^1_p \times [0, 1]$ to a once-punctured torus $\Sigma_{1,1}$
sends the annulus trace to the genus-$1$ curvature:
\begin{equation}\label{eq:thqg-annulus-to-kappa}
\Delta_{\mathrm{ns}}^{\mathrm{oc}}\!\bigl(\operatorname{Tr}_\cA\bigr)
\;=\;
\kappa(\cA) \cdot \lambda_1,
\end{equation}
where
$\lambda_1 \in H^2(\overline{\mathcal{M}}_{1,1})$ is the first
Chern class of the Hodge bundle.
The categorical annulus trace theorem supplies the source of
$\Delta_{\mathrm{ns}}^{\mathrm{oc}}$ and the class
$\operatorname{Tr}_\cA$; the scalar $\kappa(\cA)$ enters through
the closed-sector normalization above.
\end{proposition}

\begin{proof}
The non-separating degeneration identifies the two boundary
components of the annulus $S^1 \times [0, 1]$ to create a
genus-$1$ surface. The annulus trace theorem identifies the annulus
amplitude with the categorical Hochschild class
$\operatorname{Tr}_\cA=[1]$ in
$\mathrm{HH}_{0,\mathrm{cat}}(\mathcal M)$. The cyclic
Calabi--Yau class $\eta_{\mathcal M}$ is the datum that lets this
chain-level class be paired across the two boundary components.

By hypothesis, the resulting sewing operation has fixed type: its
closed projection is fixed to be
$\Theta_\cA^{(1,0)}$, the genus-$1$ component of the bar-intrinsic
MC element $\Theta_\cA=D_\cA-d_0$
(Theorem~\ref{thm:mc2-bar-intrinsic}). Theorem~D
(Theorem~\ref{thm:modular-characteristic}) gives the scalar
normalization
$\Theta_\cA^{(1,0)}=\kappa(\cA)\lambda_1$ on the modular Koszul
locus. Substitution gives
$\Delta_{\mathrm{ns}}^{\mathrm{oc}}(\operatorname{Tr}_\cA)
=\kappa(\cA)\lambda_1$.
\end{proof}

% ===================================================================
\subsection{The open/closed MC element}
\label{subsec:thqg-oc-mc-element}
\index{Maurer--Cartan element!open/closed|textbf}
% ===================================================================

\begin{definition}[Admissible open/closed extension]
\label{def:thqg-admissible-oc-extension}
Let $(\cA, X, \langle \cdot,\cdot\rangle,k)$ be a gravitational
input on a tangential log curve $(X,D,\tau)$, with boundary modules
$\mathcal{M}_1,\ldots,\mathcal{M}_r$. An \emph{admissible
open/closed extension} consists of:
\begin{enumerate}[label=\textup{(\roman*)}]
\item $A_\infty$-module structures
 $\mu^{\mathcal{M}_j}=\sum_{n\geq 1}\mu_n^{\mathcal{M}_j}$ on the
 boundary modules;
\item continuous mixed maps for each Type~III boundary face,
 \[
 \rho_j^B\colon
 \cA\otimes\mathcal{M}_j^{\otimes |B|}
 \longrightarrow
 \mathcal{M}_j,
 \]
 and their higher closed-input brace iterates, compatible with the
 local $E_2$ Swiss-cheese action of
 Definition~\ref{def:thqg-local-oc-pair};
 \item completion for the genus-degree and pole-order filtrations,
 so that the series below converges in
 $\mathfrak{g}^{\mathrm{oc}}_{\cA,\vec{\mathcal{M}}}$;
\item compatibility with the Type~II, Type~III, and nodal
 codimension-one strata of
 $\overline{C}{}^{\,\mathrm{oc}}_{k,\vec{m}}(X,D,\tau)$, using the
 boundary and mixed differentials $d_{\mathrm{bdy}}$ and
 $d_{\mathrm{mix}}$ of~\eqref{eq:Doc-seven}.
\end{enumerate}
Thus an admissible extension is additional open/closed structure
beyond the closed bar-intrinsic element~$\Theta_\cA$ and the annulus
trace.
\end{definition}

\begin{construction}[Open/closed MC element]
\label{constr:thqg-oc-mc-element}
Let $(\cA, X, \langle \cdot, \cdot \rangle, k)$ be a gravitational
input and let $(X, D, \tau)$ be a tangential log curve with
$\cA$-modules $\mathcal{M}_1, \ldots, \mathcal{M}_r$ at the
punctures. Assume an admissible open/closed extension in the sense of
Definition~\ref{def:thqg-admissible-oc-extension}.
The open-closed convolution algebra
$\mathfrak{g}^{\mathrm{oc}}_{\cA, \vec{\mathcal{M}}}$
(Definition~\ref{def:oc-modular-convolution}) carries a dg~Lie
structure with seven-component differential~\eqref{eq:Doc-seven}.
Define the \emph{open/closed MC element}
\begin{equation}\label{eq:thqg-oc-mc}
\Theta^{\mathrm{oc}}_\cA
\;:=\;
\Theta_\cA
\;+\;
\sum_{j=1}^{r}
\mu^{\mathcal{M}_j}
\;+\;
\sum_{j=1}^{r}
\rho^{\mathcal{M}_j},
\end{equation}
where $\Theta_\cA = D_\cA - d_0 \in \gAmod$ is the bar-intrinsic
MC element (Theorem~\ref{thm:mc2-bar-intrinsic}), with $d_0$ the
genus-$0$ free differential and $D_\cA$ the full bar differential,
viewed in
$\mathfrak{g}^{\mathrm{oc}}$ via the closed-sector inclusion
\eqref{eq:Doc-seven}, and
$\mu^{\mathcal{M}_j} = \sum_{n \ge 1} \mu_n^{\mathcal{M}_j}$ is the
$A_\infty$-module structure on~$\mathcal{M}_j$. The term
$\rho^{\mathcal{M}_j}$ is the completed sum of the mixed
bulk-to-boundary maps $\rho_j^B$ and their higher closed-input brace
iterates. The $\rho$-term is the mixed-sector component; the first two
summands give the closed and pure open projections.
\end{construction}

\begin{theorem}[Open/closed MC equation; \ClaimStatusProvedHere]
\label{thm:thqg-oc-mc-equation}
\index{Maurer--Cartan equation!open/closed|textbf}
For an admissible open/closed extension, the element
$\Theta^{\mathrm{oc}}_\cA$ satisfies the Maurer--Cartan equation in
$\mathfrak{g}^{\mathrm{oc}}_{\cA, \vec{\mathcal{M}}}$:
\begin{equation}\label{eq:thqg-oc-mc-equation}
D^{\mathrm{oc}}\,\Theta^{\mathrm{oc}}_\cA
+ \tfrac{1}{2}
[\Theta^{\mathrm{oc}}_\cA, \Theta^{\mathrm{oc}}_\cA]
= 0.
\end{equation}
Here $D^{\mathrm{oc}}$ is the seven-component differential of
\eqref{eq:Doc-seven}; its closed-sector summand already contains
the non-separating clutching operator $\hbar\Delta$.
This equation decomposes by sector:
\begin{enumerate}[label=\textup{(\alph*)}]
\item \emph{Pure closed sector.}
 The restriction to $\vec{m} = 0$ recovers the closed MC equation
 $D_\cA\Theta_\cA + \tfrac{1}{2}[\Theta_\cA, \Theta_\cA]
 = 0$, with the $\hbar\Delta$ term included in the five-component
 differential $D_\cA$.
\item \emph{Pure open sector.}
 The restriction to $k = 0$ recovers the $A_\infty$-module
 equations $\sum_{i+j = n} \mu_i \circ \mu_j = 0$ at each
 boundary component.
\item \emph{Mixed sector.}
 The terms with $k \ge 1$ and $|\vec{m}| \ge 1$ encode the
 bulk-to-boundary couplings: how closed-sector operators
 (elements of $\mathcal{C}^{\bullet}_{\mathrm{ch}}(\cA, \cA)$)
 act on the boundary modules~$\mathcal{M}_j$.
\item \emph{Clutching sector.}
 The $d_{\mathrm{sew}}$ and $\hbar\Delta$ components inside
 $D^{\mathrm{oc}}$ encode the modular sewing data: separating
 nodes and non-separating genus-raising nodes.
\end{enumerate}
\end{theorem}

\begin{proof}
The equation~\eqref{eq:thqg-oc-mc-equation} is the
Stokes theorem on the bordered FM compactification
$\overline{C}{}^{\,\mathrm{oc}}_{k, \vec{m}}(X, D, \tau)$
(Construction~\ref{constr:bordered-fm}). The total boundary of
$\overline{C}{}^{\,\mathrm{oc}}_{k, \vec{m}}$ decomposes into
four types (interior collision, consecutive boundary collision,
mixed bubbling, nodal clutching), and the vanishing of the total
codimension-one boundary contribution gives precisely the four-sector
decomposition above. The admissibility hypotheses identify the
consecutive boundary faces with the operations
$\mu^{\mathcal{M}_j}$, the mixed bubbling faces with the maps
$\rho^{\mathcal{M}_j}$, and the nodal faces with
$d_{\mathrm{sew}}$ and $\hbar\Delta$.

The pure closed sector (a) is the restriction to the interior FM
compactification $\operatorname{FM}_k(X \mathbin{|} D)$, which
is the known MC equation for $\Theta_\cA$. The pure open sector (b)
is the restriction to the associahedral factor
$\prod_j K_{m_j}$, giving the $A_\infty$ relations. The mixed
sector (c) comes from the mixed boundary strata (Type III faces),
and the clutching sector (d) from the nodal boundary strata
(Type IV faces).
\end{proof}

% ===================================================================
\subsection{The open/closed projection principle}
\label{subsec:thqg-oc-projection}
\index{projection principle!open/closed|textbf}
% ===================================================================

\begin{theorem}[Open/closed projection principle; \ClaimStatusProvedHere]
\label{thm:thqg-oc-projection}
The open/closed MC element contains the closed, open, mixed, and
shadow projections of the seven-entry holographic modular Koszul
datum:
\[
\mathcal{H}(\cA)
=
(\cA,\cA^i,\cA^!,\mathbf{C}_{\mathrm{ch}}(\cA),
r(z),\Theta_\cA,
\nabla^{\mathrm{hol}}).
\]
Here $\cA^i=H^*(B^{\mathrm{ch}}(\cA))$ is the bar-dual coalgebra,
whereas $\cA^!$ is the Verdier-dual line algebra on the finite-type
or completed continuous Koszul locus, and
$\mathbf{C}_{\mathrm{ch}}(\cA)
:=\mathbf{Z}^{\mathrm{der}}_{\mathrm{ch}}(\cA)$ is the
cochain-level chiral derived center, hence the intrinsic
closed-sector slot. The
line category
$\mathsf{Line}(\cA)=\cA^!\text{-}\mathsf{mod}$ is attached to
the third entry after Verdier/Koszul rectification. The projections
are:
\begin{enumerate}[label=\textup{(\roman*)}]
\item The \emph{closed projection}
 $\pi_{\mathrm{cl}} \colon
 \mathfrak{g}^{\mathrm{oc}} \to \gAmod$
 sends $\Theta^{\mathrm{oc}}_\cA \mapsto \Theta_\cA$.

\item The \emph{open projection}
 $\pi_{\mathrm{op}} \colon
 \mathfrak{g}^{\mathrm{oc}} \to
 \prod_j \operatorname{End}(\mathcal{M}_j)$
 sends $\Theta^{\mathrm{oc}}_\cA \mapsto
 \sum_j \mu^{\mathcal{M}_j}$.

\item The \emph{mixed projection} at genus~$0$, one interior
 point, and $n$~boundary points on one component gives the
 \emph{bulk-to-boundary map}
 $\iota_{0,1,n} \colon
 \mathcal{C}^{\bullet}_{\mathrm{ch}}(\cA, \cA) \to
 \operatorname{Hom}(\mathcal{M}_j^{\otimes n}, \mathcal{M}_j)$,
 which is the universal open/closed coupling
 (Theorem~\ref{thm:thqg-swiss-cheese}) realized geometrically.

\item The \emph{shadow obstruction tower projections}
 $\Theta^{\mathrm{oc}}_\cA|_{\le r}$ at each degree~$r$ recover
 the nonlinear shadows
 $(\kappa(\cA), \mathfrak{C}, \mathfrak{Q}, \ldots)$
 in the closed sector, and the \emph{open boundary corrections}
 $(\mu_2^{\mathcal{M}}, \mu_3^{\mathcal{M}}, \ldots)$
 in the open sector.
\end{enumerate}
Recovering $\cA^i$, $\cA^!$, $r(z)$, or
$\nabla^{\mathrm{hol}}$ requires their respective slots in the
completed datum. The
coalgebra $\cA^i$ is the cohomology of the chiral bar coalgebra;
$\cA^!$ is obtained only after Verdier duality on $\cA^i$ with the
finite-type or completed continuous hypotheses; $r(z)$ comes from the
binary collision residue; and $\nabla^{\mathrm{hol}}$ comes from the
modular connection data.
\end{theorem}

\begin{proof}
Each statement follows from the additive decomposition
$\Theta^{\mathrm{oc}}_\cA
= \Theta_\cA + \sum_j \mu^{\mathcal{M}_j}
 + \sum_j \rho^{\mathcal{M}_j}$
and the seven-component structure~\eqref{eq:Doc-seven} of the
differential $D^{\mathrm{oc}}$. The closed projection (i) drops
the boundary terms ($d_{\mathrm{bdy}}$ and $d_{\mathrm{mix}}$),
recovering the five-component closed differential. The open
projection (ii) drops the interior terms. The mixed projection (iii)
extracts the Type~III cross terms, whose one-interior,
one-boundary-component part lives in
$\operatorname{Hom}(C_*(\overline{C}^{\,\mathrm{oc}}_{1, (n)}),
\operatorname{Hom}(\cA \otimes \mathcal{M}_j^{\otimes n},
\mathcal{M}_j))$.
At genus~$0$ on a formal disk, this identifies with the chiral
Hochschild cochain complex by
Theorem~\ref{thm:thqg-local-global-bridge}(ii).
The shadow obstruction tower projections (iv) follow from the degree
filtration on $\mathfrak{g}^{\mathrm{oc}}$. The final paragraph is
the object separation: a projection of
$\Theta^{\mathrm{oc}}_\cA$ supplies closed, open, mixed, and filtered
MC components, while the bar-dual coalgebra, Verdier dual algebra,
collision residue, and holomorphic connection are separate
constructions with their own hypotheses.
\end{proof}

\begin{theorem}[MC-forced open-closed consistency; \ClaimStatusProvedHere]
\label{thm:thqg-mc-forced-consistency}
\index{open-closed consistency!MC-forced|textbf}
\index{Maurer--Cartan equation!open-closed consistency|textbf}
Let $\cA$ be a modular Koszul chiral algebra with boundary
category~$\mathcal{M}$, and let $\Theta^{\mathrm{oc}}_\cA
= \Theta_\cA + \sum_j \mu^{\mathcal{M}_j}
 + \sum_j \rho^{\mathcal{M}_j}$
be the open/closed MC element
\textup{(}Construction~\textup{\ref{constr:thqg-oc-mc-element})}.
The modular MC equation
\textup{(}Theorem~\textup{\ref{thm:thqg-oc-mc-equation})}
imposes a compatible system at every genus~$g \geq 0$:
\begin{enumerate}[label=\textup{(\alph*)}]
\item \emph{Closed-sector genus tower.}
 The closed projection
 $\pi_{\mathrm{cl}}(\Theta^{\mathrm{oc}}_\cA)|_{g}
 = \Theta_\cA|_{g}$
 recovers the genus-$g$ component of the bar-intrinsic MC element.
On the uniform-weight lane, this is the genus tower
$F_g = \kappa(\cA) \cdot \lambda_g$
\textup{(}Theorem~\textup{\ref{thm:mc2-bar-intrinsic}}(ii),
 Theorem~D\textup{)}.

\item \emph{Open-sector module structure.}
 The open projection
 $\pi_{\mathrm{op}}(\Theta^{\mathrm{oc}}_\cA)|_{g}$
 recovers the $A_\infty$-module operations
 $\mu^{\mathcal{M}_j}_{n}$ on the boundary components.

\item \emph{Mixed bulk-to-boundary sector.}
 The mixed projection recovers the maps
 $\rho^{\mathcal{M}_j}$ of
 Definition~\ref{def:thqg-admissible-oc-extension}. If a cyclic
 boundary family is equipped with a mixed sewing map compatible with
 $\Theta_\cA|_g$, then the mixed differential inserts the image of
 the closed curvature class in the open module equation. This
 curvature insertion is open/closed MC data supplied in addition to
 the categorical Hochschild trace.

\item \emph{Clutching sector.}
 The $\hbar\Delta$ summand of $D^{\mathrm{oc}}$ encodes the
 genus-raising sewing: non-separating clutching
 $\xi_{\mathrm{ns}}^*(\Theta^{(g+1)})
 = \Delta_{\mathrm{ns}}(\Theta^{(g)})$
 and separating clutching
 $\xi_{\mathrm{sep}}^*(\Theta^{(g)})
 = \sum_{g_1+g_2=g}
 \Theta^{(g_1)} \star \Theta^{(g_2)}$
 \textup{(}Theorem~\textup{\ref{thm:mc2-bar-intrinsic}}(iii)\textup{)}.
\end{enumerate}
Since the closed, open, mixed, and clutching sectors are projections
of the same
MC element~$\Theta^{\mathrm{oc}}_\cA$, the genus-$g$ component
of the MC equation is a single identity whose restriction to each
sector yields a compatible system:
\begin{enumerate}[resume, label=\textup{(\alph*)}]
\item \emph{Open-closed consistency at every genus.}
	 The closed-sector genus-$g$ component and the open-sector
	 boundary operations are projections of the chosen element
	 $\Theta^{\mathrm{oc}}_\cA|_g$. The MC equation constrains their
	 joint values in the open/closed convolution algebra. Reconstructing
	 the closed scalar tower from
	 $\mathrm{HH}_{\bullet,\mathrm{cat}}(\mathcal M)$, or reconstructing
	 the annulus trace from Theorem~D, requires additional comparison
	 data.
\end{enumerate}
\end{theorem}

\begin{proof}
The MC equation
$D^{\mathrm{oc}}\Theta^{\mathrm{oc}} + \tfrac{1}{2}
[\Theta^{\mathrm{oc}}, \Theta^{\mathrm{oc}}] = 0$
holds at all genera by
Theorem~\ref{thm:thqg-oc-mc-equation}. Each genus-$g$
component of this equation lives in
$\mathfrak{g}^{\mathrm{oc},(g)}_{\cA, \vec{\mathcal{M}}}$,
which decomposes into the four sectors of
Theorem~\ref{thm:thqg-oc-mc-equation}(a)--(d).
The projections (a)--(d) follow from
Theorem~\ref{thm:thqg-oc-projection}: the closed projection
recovers $\Theta_\cA|_g$ with the scalar trace
$\kappa(\cA)\lambda_g$ on the uniform-weight lane;
the open projection recovers the module operations; the mixed
projection recovers the bulk-to-boundary maps and only inserts
closed curvature terms through the specified mixed sewing map;
the clutching projection recovers the separating and
non-separating degenerations.

Part~(e) is the structural consequence. The genus-$g$
MC equation is a \emph{single} vanishing condition in
$\mathfrak{g}^{\mathrm{oc},(g)}$. Its restriction to any
sector is a consequence of the same equation. Therefore
any identity obtained by projecting to the closed sector
is compatible with the identity obtained by projecting to the
open sector. Compatibility means that both identities are projections
of the same source equation. The
compatibilities are projections of the Stokes theorem on the
bordered FM compactification
$\overline{C}{}^{\,\mathrm{oc}}_{k, \vec{m}}$.
\end{proof}

\begin{remark}[Scalar and open-sector clauses]
\label{rem:thqg-mc-forced-scope}
Theorem~\ref{thm:thqg-mc-forced-consistency} proves compatibility
of the chosen open, closed, mixed, and clutching projections at
every genus. Two further clauses have independent hypotheses:
\begin{enumerate}[label=\textup{(\roman*)}]
\item \emph{The scalar identification.}
The specific formula $F_g = \kappa(\cA) \cdot \lambda_g$
is proved on the uniform-weight lane and at genus~$1$ for all
 families. For multi-weight algebras at genus~$g \geq 2$,
 the scalar formula \emph{fails}: the free energy receives a
 cross-channel correction $\delta F_g^{\mathrm{cross}} \neq 0$
 \textup{(}Theorem~\ref{thm:multi-weight-genus-expansion}\textup{)}.
 The scalar trace of~$\Theta_\cA$ is therefore a tautological-class
 calculation on $\overline{\cM}_g$, separate from open/closed
 compatibility.
\item \emph{The independent open-sector derivation.}
 Computing $F_g$ from iterated clutching on
 $\mathcal{C}_{\mathrm{op}}$ alone (without invoking the
 bar-intrinsic construction) and verifying agreement with the
 closed-sector tower is a separate construction. The MC-forced
 consistency gives the compatibility target for that calculation.
\end{enumerate}
\end{remark}

% ===================================================================
\subsection{CT-2: the modular cooperad conjecture for the open sector}
\label{subsec:thqg-ct2-cooperad-conjecture}
\index{CT-2 conjecture|textbf}
\index{modular cooperad!open sector|textbf}
% ===================================================================

The open-sector factorization dg category $\mathcal{C}_{\mathrm{op}}$
carries cocomposition maps read off from clutching on bordered
Fulton--MacPherson compactifications, and at genus~$1$ the annulus
trace gives the categorical Hochschild chain object of
Theorem~\ref{thm:thqg-annulus-trace}. The modular cooperad
conjecture asks for two additional structures: assembly of these
chain-level cocompositions into a Getzler--Kapranov modular
cooperad \cite{GetzlerKapranov98}, and compatibility of that
cooperad with the open-closed MC element
$\Theta^{\mathrm{oc}}_\cA$. The categorical trace is the genus-$1$
open input for this cooperad clause.

\begin{conjecture}[CT-2: modular cooperad on the open sector]
\ClaimStatusConjectured
\label{conj:ct2}
Let $\cA$ be a gravitational input
(Definition~\ref{def:thqg-standing-hypotheses}) and let
$\mathcal{C}_{\mathrm{op}}$ denote the open-sector factorization
dg category of
\S\ref{subsec:thqg-open-sector-factorization-category}.
There exists a modular cooperad structure
\[
\Delta^{\mathrm{op}}_{\mathrm{mod}}
\colon
\mathcal{C}_{\mathrm{op}}
\longrightarrow
\mathcal{C}_{\mathrm{op}} \mathbin{\widehat{\boxtimes}}
\mathcal{C}_{\mathrm{op}}
\]
on $\mathcal{C}_{\mathrm{op}}$ with the following three properties.
\begin{enumerate}[label=\textup{(\alph*)}]
\item \emph{Clutching compatibility.}
 The cocomposition $\Delta^{\mathrm{op}}_{\mathrm{mod}}$ is compatible
 with the clutching maps on the bordered Fulton--MacPherson
 compactification $\overline{C}{}^{\,\mathrm{oc}}_{k,\vec{m}}$
 of \S\ref{sec:bordered-fm}.
\item \emph{Annulus trace at genus~$1$.}
 The trace of $\Delta^{\mathrm{op}}_{\mathrm{mod}}$ at genus~$1$
 recovers the annulus identification of
 Theorem~\ref{thm:thqg-annulus-trace}.
\item \emph{MC determination.}
 The full structure is determined by the open-closed MC element
 $\Theta^{\mathrm{oc}}_\cA$ of
 Definition~\ref{constr:thqg-oc-mc-element}: two modular cooperad
 structures on $\mathcal{C}_{\mathrm{op}}$ that agree on
 $\Theta^{\mathrm{oc}}_\cA$ are equivalent as
 $(\infty,2)$-functors.
\end{enumerate}
\end{conjecture}

\begin{remark}[Existence and coherence obligations]
\label{rem:ct2-existence-coherence-obligations}
Conjecture~\ref{conj:ct2} has three logically separate clauses.

\emph{Existence.}
An existence proof must lift the known modular operad structure on
the closed sector $\mathcal{C}_{\mathrm{cl}}$ through the bar-cobar
adjunction of Theorem~A. The required inputs are homotopy Koszulness
for the open-sector category from the Vol.~II PVA descent and
compatibility of the resulting homotopies with bordered clutching.

\emph{Clutching compatibility.}
The Getzler--Kapranov modular operad framework
\cite{GetzlerKapranov98} gives the target coherence identities.
The clause requires the cocompositions read from the bordered
Fulton--MacPherson compactification to satisfy those identities in
the open factorization category.

\emph{$(\infty,2)$-uniqueness.}
Uniqueness of the cooperad structure requires integrating the
strict clutching data to an $(\infty,2)$-category. The categorical
obstruction lies in that integration. Brochier--Woike's construction for
factorization homology and topological field theory
\cite{BrochierWoike2022}, which produces an $(\infty,2)$-functor
from clutching data under additional coherence hypotheses, is the
model for this integration clause.
\end{remark}

\begin{remark}[Inputs and obstruction classes]
\label{rem:ct2-proved-conjectural-loci}
The established inputs for Conjecture~\ref{conj:ct2} are:
\begin{enumerate}[label=\textup{(\roman*)}]
\item The bordered Fulton--MacPherson compactification
 $\overline{C}{}^{\,\mathrm{oc}}_{k,\vec{m}}(X, D, \tau)$ with its
 clutching and gluing strata is constructed in the Vol~II chapter
 on bordered FM geometry, building on
 \cite{KontsevichSoibelman}; see
 \S\ref{sec:bordered-fm}.
\item The open-closed differential squares to zero on
 the chain complex, with the convolution-level identity proved via
 Mok's log-FM framework and the ambient differential identity
 recorded in Theorem~\ref{thm:thqg-oc-mc-equation}(a). This fixes
 the ambient cochain model on which any cooperad structure must act.
\item The annulus trace theorem identifies the genus-$1$ open trace:
 \[
 \int_{S^1_p}\mathcal{C}_{\mathrm{op}}
 \simeq
 \mathrm{HH}_{\bullet,\mathrm{cat}}(\mathcal{C}_{\mathrm{op}}).
 \]
 This fixes the chain object that Conjecture~\ref{conj:ct2}(b)
 must recover. The trace of
 $\Delta^{\mathrm{op}}_{\mathrm{mod}}$ equals the genus-$1$
 component of $\Theta^{\mathrm{oc}}_\cA$ exactly in the compatibility
 clause of the conjecture.
\end{enumerate}
The obstruction classes are:
\begin{enumerate}[label=\textup{(\roman*)}]
\item the lift of the modular cooperad structure along
 the bar-cobar adjunction from
 $\mathcal{C}_{\mathrm{cl}}$ to $\mathcal{C}_{\mathrm{op}}$,
 using homotopy Koszulness of $\mathcal{C}_{\mathrm{op}}$ (Vol~II
 PVA descent) to promote the strict clutching data to a cooperad
 up to coherent homotopy, without identifying the categorical
 trace with chiral Hochschild cochains except through the comparison
 maps of Theorem~\ref{thm:thqg-hochschild-register-separation}.
\item the integration of the homotopy-coherent clutching
 data to an $(\infty,2)$-category together with a proof that two such
 integrations agreeing on $\Theta^{\mathrm{oc}}_\cA$ are equivalent
 as $(\infty,2)$-functors, following the Brochier--Woike
 localization mechanism \cite{BrochierWoike2022}.
\end{enumerate}
These two classes are the existence and uniqueness content of
Conjecture~\ref{conj:ct2}.
\end{remark}

\begin{remark}[Categorical localization obstruction]
\label{rem:ct2-brochier-woike}
\index{Brochier--Woike localization|textbf}
The obstruction is categorical. The bordered FM geometry, the
chain-level differential, and the genus-$1$ annulus trace are
constructed at the level of dg categories and their homologies;
clutching compatibility is a concrete chain-level identity. The
classification of cooperad structures up to equivalence requires a
homotopy-coherent $(\infty,2)$-category of open-sector
factorization algebras receiving these strict chain-level data.

Brochier--Woike \cite{BrochierWoike2022} first builds a
$(2,2)$-category of bordered
factorization data with strict clutching, and then localizes at
the homotopy equivalences using the universal property of
$(\infty,2)$-localization. The categorical condition is that
the localization preserves the clutching compatibility recorded in
Conjecture~\ref{conj:ct2}(a). Under this condition, the
Brochier--Woike rigidity argument applied to
$\Theta^{\mathrm{oc}}_\cA$ gives the uniqueness statement of
Conjecture~\ref{conj:ct2}(c). The obstruction is the preservation
of clutching under $(\infty,2)$-localization.
\end{remark}

% ===================================================================
\subsection{The first nonlinear open/closed invariant}
\label{subsec:thqg-nonlinear-oc-invariant}
\index{open/closed invariant!nonlinear|textbf}
% ===================================================================

\begin{construction}[Open/closed quartic resonance class]
\label{constr:thqg-oc-quartic-resonance}
Let $\cA$ be a gravitational input of shadow archetype M
(Definition~\ref{def:thqg-shadow-archetype}, mixed class with
$r_{\max} = \infty$). The quartic contact invariant
$\mathfrak{Q}^{\mathrm{ct}}(\cA) \in
\mathcal{C}^2_{\mathrm{ch}}(\cA, \cA)$
is a degree-$2$ chiral Hochschild cochain measuring the quartic
nonlinearity of the OPE. Define the \emph{open/closed quartic
resonance class}
\begin{equation}\label{eq:thqg-oc-quartic}
\mathfrak{R}^{\mathrm{oc}}_{4}(\cA)
\;:=\;
[\mathfrak{Q}^{\mathrm{ct}}(\cA)]
\;+\;
[\mathfrak{C}(\cA) \star_{P_\cA} \mathfrak{C}(\cA)]
\;\in\;
H^2\!\bigl(\mathcal{C}^{\bullet}_{\mathrm{ch}}(\cA, \cA),\, \delta\bigr)
\;=\; \mathcal{Z}^2,
\end{equation}
where $\mathfrak{C}(\cA)$ is the cubic shadow
(eq.~\eqref{eq:nms-cubic-shadow}),
$H_\cA$ is the Hessian (the quadratic shadow, i.e.\ the modular
characteristic $\kappa(\cA)$ on the one-dimensional primary slice),
$P_\cA = H_\cA^{-1}$ is the propagator, and
$(\mathfrak{C} \star_{P_\cA} \mathfrak{C})(x_1, x_2, x_3, x_4)
:= \sum_\alpha \mathfrak{C}(x_1, x_2, e_\alpha)\, P_\cA^{\alpha\beta}\,
\mathfrak{C}(e_\beta, x_3, x_4)$
is the sewing product contracting one pair of indices with the
propagator. The first term is the local quartic contact;
the second is the tree correction from two cubic emissions joined
by one propagator.
\end{construction}

\begin{theorem}[Quartic resonance on the four shadow witnesses;
\ClaimStatusProvedHere]
\label{thm:thqg-oc-quartic-vanishing}
\begin{enumerate}[label=\textup{(\roman*)}]
\item \emph{Gaussian vanishing.}
 For $\cA$ of shadow archetype~G (Heisenberg, lattice, free fermion):
 $\mathfrak{R}^{\mathrm{oc}}_4(\cA) = 0$.

\item \emph{Lie/tree vanishing of contact.}
 For $\cA$ of archetype~L (affine $\widehat{\mathfrak{g}}$ at
 generic level): the contact term vanishes in the standard gauge,
 $[\mathfrak{Q}^{\mathrm{ct}}(\cA)] = 0$, but the tree correction
 $[\mathfrak{C} \star_P \mathfrak{C}]$ is generically nonzero.

\item \emph{Contact termination.}
 For $\cA$ of archetype~C
 \textup{(}the standard conformal-weight family
 $\beta\gamma_\lambda$ or $bc_\lambda$\textup{)}, the cubic shadow
 vanishes, the quartic contact term is the last nonzero correction,
 and the quintic collapses by rank-one abelian rigidity on the
 internal weight-changing slice
 \textup{(}Corollary~\ref{cor:nms-betagamma-mu-vanishing};
 Proposition~\ref{prop:depth-gap-trichotomy}\textup{)}.

\item \emph{Virasoro mixed nonvanishing.}
 For $\cA=\mathrm{Vir}_c$ on the non-singular shadow surface
 $c(5c+22)\ne0$, both terms are generically nonzero, and
 \begin{equation}\label{eq:thqg-oc-quartic-virasoro}
 \mathfrak{R}^{\mathrm{oc}}_{4}(\mathrm{Vir}_c)
 =
 \left[\frac{10}{c(5c + 22)}\right] \cdot x^4
 + [\mathfrak{C}_{\mathrm{Vir}} \star_{P_{\mathrm{Vir}}}
 \mathfrak{C}_{\mathrm{Vir}}]
 \;\neq\; 0
 \end{equation}
 for generic~$c$. The contact term has poles at $c = 0$ and
 $c = -22/5$ (the null-state divisor), while the tree term is
 regular.
\end{enumerate}
For a principal $\mathcal{W}_N$ algebra, the same class-M
nonvanishing statement holds on each generic component where the
landscape-census quartic shadow coefficient in the chosen channel is
nonzero.
\end{theorem}

\begin{proof}
(i)--(iii) follow directly from the shadow archetype classification
and the explicit computations in the nonlinear modular shadows
appendix (Appendix~\ref{app:nonlinear-modular-shadows}):
Heisenberg has $\mathfrak{C} = \mathfrak{Q} = 0$;
affine has $\mathfrak{Q}^{\mathrm{ct}} = 0$ (Jacobi identity kills
the quartic obstruction) but
$\mathfrak{C}(x, y, z) = \kappa(x, [y, z]) \neq 0$;
$\beta\gamma_\lambda$ has $\mathfrak{C} = 0$ on the standard
conformal-weight family, with the quartic contact term as the last
nonzero correction; the internal weight-changing line satisfies
$\mu_{\beta\gamma} = 0$ by rank-one abelian rigidity.

(iv) For Virasoro: $\mathfrak{C}_{\mathrm{Vir}} = 2x^3$ from
the universal gravitational cubic tensor, and
$\mathfrak{Q}^{\mathrm{ct}}_{\mathrm{Vir}}
= 10 / (c(5c + 22))$ from the quartic contact computation.
Both are nonzero for generic~$c$, and they contribute independently
to $\mathfrak{R}^{\mathrm{oc}}_4$ since contact lives in the quartic
sector of $\mathcal{C}^2_{\mathrm{ch}}$ while the tree term lives
in the image of the sewing product
$\mathcal{C}^1_{\mathrm{ch}} \star \mathcal{C}^1_{\mathrm{ch}}
\to \mathcal{C}^2_{\mathrm{ch}}$.
The principal $\mathcal{W}_N$ assertion is the same computation
with the channel coefficient supplied by
Chapter~\ref{ch:landscape-census}; nonzero quartic shadow
coefficient is precisely the genericity hypothesis in the statement.
\end{proof}

\begin{remark}[The open/closed resonance class as gravitational observable]
\label{rem:thqg-oc-gravitational-observable}
\index{gravitational observable|textbf}
The class $\mathfrak{R}^{\mathrm{oc}}_4(\cA) \in \mathcal{Z}^2$
is the first genuinely \emph{nonlinear} gravitational observable
in the holographic modular Koszul datum. The scalar modular
characteristic $\kappa(\cA)$ and its duality constraint (Theorem~D)
are linear invariants of the boundary OPE. The cubic shadow
$\mathfrak{C}(\cA)$ is the first nonlinear correction but lives
in $\mathcal{Z}^1$ (derivations), not in $\mathcal{Z}^2$
(obstructions). The quartic class $\mathfrak{R}^{\mathrm{oc}}_4$
is the first obstruction-level nonlinear invariant, and its
nonvanishing for $\mathcal{W}$-type theories signals that the
boundary $A_\infty$ structure is \emph{irreducibly quartic}:
quartic terms are necessary in the modular master equation for the
perturbative partition function.

The clutching law for $\mathfrak{R}^{\mathrm{oc}}_4$
(Theorem~\ref{thm:nms-clutching-law-modular-resonance}) shows
that this obstruction propagates across nodal degenerations,
and the tree correction $\mathfrak{C} \star_P \mathfrak{C}$ is
the cochain trace of cubic exchange across a boundary node. Under a
physical bulk comparison of the form
Proposition~\ref{prop:thqg-universal-action-not-reconstruction},
the same class is read as the quartic bulk channel: a local
four-point contact contribution
$\mathfrak{Q}^{\mathrm{ct}}$ plus a tree contribution
$\mathfrak{C} \star_P \mathfrak{C}$. Without that comparison, it is
a class in the chiral Hochschild cochain center.
\end{remark}

\begin{remark}[$\mathcal{W}_3$ and the first nonlinear composite]
\label{rem:thqg-w3-nonlinear}
\index{W3 algebra@$\mathcal{W}_3$ algebra!nonlinear composite|textbf}
The $\mathcal{W}_3$ algebra is the minimal example where
$\mathfrak{R}^{\mathrm{oc}}_4$ detects genuinely nonlinear
composite structure. The $W$--$W$ OPE produces the quadratic
composite $\Lambda := {:}TT{:} - \tfrac{3}{10}\partial^2 T$,
which lies outside the linear span of the strong generators.
The quartic contact coefficient receives contributions from
both $T$--$T$--$T$--$T$ and $T$--$T$--$W$--$W$ channels.
The quartic resonance divisor for $\mathcal{W}_3$ contains the
factor $c(2c - 1)(5c + 22)$
(cf.\ the discussion in
Appendix~\ref{app:nonlinear-modular-shadows}),
where $c = 1/2$ is the Ising model (minimal model $M(4, 3)$)
and $c = -22/5$ is the Yang--Lee edge singularity
(minimal model $M(2, 5)$). These are the first two points
where the boundary shadow obstruction tower degenerates.
\end{remark}

% ===================================================================
\subsection{The modular Koszul datum, completed}
\label{subsec:thqg-platonic-completed}
\index{platonic datum!completed|textbf}
% ===================================================================

\begin{definition}[Completed modular Koszul datum]
\label{def:thqg-completed-platonic-datum}
Let $(\cA, X, \langle \cdot, \cdot \rangle, k)$ be a gravitational
input on a tangential log curve $(X, D, \tau)$ with boundary
modules $\vec{\mathcal{M}}$. The \emph{completed modular Koszul datum}
is the nonuple
\begin{equation}\label{eq:thqg-completed-platonic}
\Pi^{\mathrm{oc}}_X(\cA)
\;:=\;
\Bigl(
\cA,\;\;
\cA^i,\;\;
\cA^!,\;\;
\mathbf{Z}^{\mathrm{der}}_{\mathrm{ch}}(\cA),\;\;
\Theta^{\mathrm{oc}}_\cA,\;\;
r(z),\;\;
\nabla^{\mathrm{hol}},\;\;
\operatorname{Tr}_\cA,\;\;
\mathfrak{R}^{\mathrm{oc}}_{\bullet}(\cA)
\Bigr)
\end{equation}
consisting of:
\begin{center}
\renewcommand{\arraystretch}{1.3}
\small
\begin{tabular}{cll}
\textbf{Component} & \textbf{Object} & \textbf{Source} \\
\hline
$\cA$ & boundary chiral algebra & given \\
$\cA^i$ & bar-dual coalgebra & $H^*(B^{\mathrm{ch}}(\cA))$ \\
$\cA^!$ & Verdier-dual line algebra & Verdier dual of $\cA^i$ \\
$\mathbf{Z}^{\mathrm{der}}_{\mathrm{ch}}(\cA)$
 & cochain-level derived center (universal closed-sector actor)
 & Def.~\ref{def:thqg-chiral-derived-center} \\
$\Theta^{\mathrm{oc}}_\cA$
 & open/closed MC element & Constr.~\ref{constr:thqg-oc-mc-element} \\
$r(z)$ & collision residue (classical $r$-matrix)
 & Def.~\ref{def:holographic-modular-koszul-datum} \\
$\nabla^{\mathrm{hol}}$
 & modular shadow connection & Def.~\ref{def:holographic-modular-koszul-datum} \\
$\operatorname{Tr}_\cA$
 & annulus trace (categorical $\mathrm{HH}_\bullet$)
 & Thm.~\ref{thm:thqg-annulus-trace} \\
$\mathfrak{R}^{\mathrm{oc}}_\bullet(\cA)$
 & nonlinear resonance tower & Constr.~\ref{constr:thqg-oc-quartic-resonance}
\end{tabular}
\end{center}
The datum extends the holographic modular Koszul datum
$\mathcal{H}(\cA)$
(Definition~\ref{def:holographic-modular-koszul-datum})
by writing the fourth slot
$\mathbf{C}_{\mathrm{ch}}(\cA)$ as its cochain model
$\mathbf{Z}^{\mathrm{der}}_{\mathrm{ch}}(\cA)$, by replacing
$\Theta_\cA$ with its open/closed lift
$\Theta^{\mathrm{oc}}_\cA$, and by adding the annulus trace and the
nonlinear resonance tower. The entries $\cA^i$, $\cA^!$, and
$\mathbf{Z}^{\mathrm{der}}_{\mathrm{ch}}(\cA)$ are distinct objects:
$\cA^i$ is coalgebraic bar cohomology, $\cA^!$ is Verdier/Koszul dual
after finite-type or completed continuous duality, and
$\mathbf{Z}^{\mathrm{der}}_{\mathrm{ch}}(\cA)$ is chiral Hochschild
cochains, the closed-sector entry of the seven-entry datum.
Here $\operatorname{Tr}_\cA$ means
$\int_{S^1}\mathcal M\simeq
\mathrm{HH}_{\bullet,\mathrm{cat}}(\mathcal M)$; the chiral cochain
model $\mathcal{C}^{\bullet}_{\mathrm{ch}}(\cA,\cA)$ is the separate
fourth slot.
\end{definition}

\begin{remark}[What the completed datum captures]
\label{rem:thqg-completed-datum-captures}
The completed datum $\Pi^{\mathrm{oc}}_X(\cA)$ packages the
following typed mathematical and physical content:
\begin{enumerate}[label=\textup{(\alph*)}]
\item \emph{Boundary.} $\cA$ is the boundary chiral algebra
 (the ``boundary CFT'' of holography).
\item \emph{Bulk.}
 $\mathbf{Z}^{\mathrm{der}}_{\mathrm{ch}}(\cA)$ is the
 universal cochain closed-sector object functorially attached to
 the boundary. Its cohomology
 $\mathcal{Z}^{\mathrm{der}}_{\mathrm{ch}}(\cA)$ is the on-shell
 Gerstenhaber algebra of bulk fields after a physical comparison map
 has been supplied. The perturbative gravitational partition
function also uses
$\Theta^{\mathrm{oc}}_\cA$, the scalar modular package, and the
clutching maps, in addition to the categorical annulus trace and
$\operatorname{THH}$ comparison slots.
\item \emph{Lines.}
 $\cA^!$ generates the line-operator category
 $\mathsf{Line}(\cA):=\cA^!\text{-}\mathsf{mod}$ when the
 Verdier--Koszul comparison supplies that category. This
 line-operator category is separate from the chiral derived center.
\item \emph{Scattering.}
 $r(z)=\operatorname{Res}^{\mathrm{coll}}_{0,2}(\Theta_\cA)$ is the
 binary genus-zero collision residue. For affine Kac--Moody,
 $r^{\mathrm{coll}}(z)=k\Omega_{\mathrm{tr}}/z$ in trace-form
 normalization, while the KZ comparison writes the same monodromy
 datum as
 $r^{\mathrm{KZ}}(z)=\Omega/((k+h^\vee)z)$ after
 Casimir and coupling rescaling. There is a single collision
 coordinate in this denominator. For Virasoro,
 $r^{\mathrm{Vir}}(z)=(c/2)z^{-3}+2Tz^{-1}$.
 A Yangian coproduct or RTT presentation requires the separate
 line-category data of
 Proposition~\ref{prop:thqg-swiss-cheese-no-yangian-coproduct}.
\item \emph{Modular completion.}
 $\Theta^{\mathrm{oc}}_\cA$ is the open/closed MC element whose
 genus and boundary projections give the shadow obstruction tower,
 boundary corrections, and clutching identities.
\item \emph{Boundary nonlinearity.}
 $\mathfrak{R}^{\mathrm{oc}}_\bullet(\cA)$ detects when the
 boundary $A_\infty$ structure becomes irreducibly nonlinear, and
 its vanishing/nonvanishing separates the shadow archetypes
 G/L/C/M. The chiral bulk is cochain-level. The topological sector
 enters through categorical trace or $\operatorname{THH}$ only after
 the relevant spectrification or Calabi--Yau comparison has been
 specified.
\end{enumerate}
\end{remark}

\section{K3 class-\texorpdfstring{$\mathcal S$}{S} comparison in the open-closed register}
\label{sec:thqg-occ-supplement}

\begin{remark}[Class-$\mathcal S$ on $\Sigma_{0,24}$ as a protected chart]
\label{rem:thqg-occ-1}
The class-$\mathcal S$ theory $\mathcal T[A_1,\Sigma_{0,24}]$ has
$c_{4d} = 107/6$, $c_{2d} = -214$, Coulomb rank $21$, flavour
$\mathfrak{su}(2)^{24}$. It supplies the protected class-$\mathcal S$
chart in the K3 comparison. The open/closed bulk identification
requires the boundary-category and bulk-comparison data of
Theorem~\ref{thm:thqg-local-global-bridge}. The proved numerical and
automorphic bridge is the composite
\[
 \mathcal I_{\mathrm{Schur}}
 \xrightarrow{\;M_{24}\text{-avg}\;}
 \phi_{0,1}^{K3}
 \xrightarrow{\;\mathrm{Borch}\;}
 \Delta_5
\]
of Vol.~III Theorem~\ref{V3-thm:schur-to-delta5-composite}, rather
than a direct equality with $1/\Delta_5$. The $24$ punctures on
$\Sigma_{0,24}$ host the $24$ flavour $\mathfrak{su}(2)$ factors and
match the $24$ Kodaira $I_1$ fibres only after the bi-based Ran datum
of Vol.~III Theorem~\ref{V3-thm:k3-bi-based-factorization} is chosen.
This is the class-$\mathcal S$ face of the K3 chiral-bialgebra
comparison \cite{Gaiotto2009,BeemRastelli2015}.
\end{remark}

\begin{remark}[Bi-based Ran factorisation and open/closed input]
\label{rem:thqg-occ-2}
The bi-based Ran space $(\mathrm{Ran}(E^{\mathrm{nod,sm}}_{24}),\,\bar{\mathcal A}_2)$
is a factorisation substrate. TFT structure enters after the boundary
category, open/closed coupling, and sewing constraints are chosen.
The chiral base
$\mathrm{Ran}(E^{\mathrm{nod,sm}}_{24})$ carries the boundary
configuration data at the Kodaira $I_1$ fibres; the parameter base
$\bar{\mathcal A}_2$ carries the genus-$2$ Siegel/Kuga--Satake
automorphic datum. Open/closed language enters only after a boundary
category $\mathcal M$, compact-generator chart, and mixed MC coupling
have been specified as in
Theorem~\ref{thm:thqg-local-global-bridge} and
Construction~\ref{constr:thqg-oc-mc-element}. The Steiner
$S(5,8,24)$ rigidity fixes the $M_{24}$-equivariant labelling of the
$24$ nodes/punctures \cite{ConwaySloane1988}; the Cardy condition is a
separate open/closed sewing constraint, independent of the Steiner
design and Kuga--Satake data. The geometric input is
Vol.~III Theorem~\ref{V3-thm:k3-bi-based-factorization}.
\end{remark}

\begin{remark}[Nekrasov $\Omega$-background on the bi-based datum;
self-dual phase $\epsilon_1 = -\epsilon_2$]
\label{rem:thqg-occ-omega-background}
The bi-based Ran datum
$(\mathrm{Ran}(E^{\mathrm{nod,sm}}_{24}),\,\bar{\mathcal A}_2)$
admits an $\Omega$-deformed comparison target. On the chiral base
$\mathrm{Ran}(E^{\mathrm{nod,sm}}_{24})$ the two equivariant
parameters $(\epsilon_1,\epsilon_2)$ act by rotation of the two
transverse planes at each Kodaira $I_1$ node. If a
holomorphic--topological prefactorisation model
$T_{\mathbf{H}_{\Delta_5}}$ lies in the boundary-linear or tempered
algebraic-sector scope of the Vol.~II
$\mathsf{SC}^{\mathrm{ch,top}}$-brace bulk theorem
\textup{(}identifying $\partial = \cA$, bulk
$= Z^{\mathrm{der}}_{\mathrm{ch}}(\cA)$, interaction
$= \mathsf{SC}^{\mathrm{ch,top}}$-brace; the gravity-line completion
to a compact dynamical-metric path integral remains conditional on
the BTZ / Cardy / Maloney--Witten / Farey-tail hypothesis package of
Remark~\ref{rem:thqg-oc-algebraic-sector}\textup{)},
then the required bulk comparison is the pair of maps
\[
\mathrm{Obs}^{\mathrm{bulk}}(T_{\mathbf{H}_{\Delta_5}})
\longrightarrow
\operatorname{Loc}_{\partial}
\mathrm{Obs}^{\mathrm{bulk}}(T_{\mathbf{H}_{\Delta_5}})
\xrightarrow{\;\beta_{T_{\mathbf H}}\;}
\mathbf{Z}^{\mathrm{der}}_{\mathrm{ch}}(\mathbf{H}_{\Delta_5})
\]
where the second arrow is an $E_2$/brace quasi-isomorphism in the
ambient completion dictated by that theorem. An equivalence with the
full $E_3$-topological physical bulk is a stronger hypothesis. The
comparison uses boundary localization, brace action, and Vol.~II
holography data in addition to the annulus trace, categorical
Hochschild homology, and $\operatorname{THH}$. With the comparison
supplied, the induced
$T^2 = U(1)_{\epsilon_1}\times U(1)_{\epsilon_2}$ action on the
derived centre is the cochain-level closed-sector side of the Nekrasov--Okounkov
equivariant trace \cite{Nekrasov2003,NekrasovOkounkov2006}. The
self-dual phase uses the independent Siegel character identity of
Vol.~III Theorem~\ref{V3-thm:k3-siegel-character}:
\[
 Z^{K3}_{\mathrm{inst}}(\epsilon_1,\epsilon_2)\big|_{\epsilon_1=-\epsilon_2=\epsilon}
 \;=\;
 \frac{1}{\Phi_{10}^{\mathrm{un}}(Z)},
 \qquad
 \hbar^2 \;=\; \epsilon_1\,\epsilon_2 \;=\; -\epsilon^2.
\]
The specialization $\hbar^2 = -1/8$ fixes
$\epsilon^2 = 1/8$, i.e.\ $\epsilon = \pm 1/(2\sqrt{2})$; this is
the self-dual $\Omega$-background point $\epsilon_1=-\epsilon_2$.
The Nekrasov--Shatashvili Bethe-Ansatz limit is the different
degeneration with one equivariant parameter sent to~$0$
\cite{NekrasovShatashvili2009}. At the self-dual point, the
Maulik--Okounkov stable-envelope $R$-matrix compares with the
$R_{\mathrm{Sieg,dyn}}$ of the K3 Hall--Drinfeld double after the
Hall--Borcherds identification and the $\hbar^2=-1/8$ specialization
are fixed \cite{MaulikOkounkov12,SchiffmannVasserot13}. The Humbert
monodromy identity $\mathrm{ord}(H_1) = K^\kappa = 2c_+ = 8$
is the matching normalization for that braid-group representation.
Primary-literature frame: Nekrasov~2003
(\texttt{hep-th/0206161}), Nekrasov--Okounkov 2006, Maulik--Okounkov
Ast\'erisque 408 (2019); AGT side
\cite{AGT09,AldayTachikawa10,CordovaGaiotto2016}.
\end{remark}

\begin{remark}[Character matching: Nekrasov partition function and chiral character]
\label{rem:thqg-occ-character-match}
For $\mathbf{H}_{\Delta_5}$ the Schur-limit of the
$\Omega$-deformed K3 partition function is computed by the
plethystic-exponential expansion of the class-$\mathcal S$
$\mathcal{T}[A_1,\Sigma_{0,24}]$ $4d$/$2d$ shadow:
\[
 \mathcal I_{\mathrm{Schur}}(q;\mathbf 1)
 \;=\;
 \mathrm{PE}\!\left[\frac{72 q - 22 q^2}{1 - q}\right]
 \;=\;
 1 + 72 q + 2678\, q^2 + 68474\, q^3 + 1351775\, q^4 + \ldots
\]
matched through $O(q^{10})$ (Beem--Lemos--Peelaers--Rastelli
JHEP~05 (2015) 020; Gadde--Rastelli--Razamat--Yan
\texttt{arXiv:1110.3740}). The $M_{24}$-average of the same series
over the $24$ $\mathfrak{su}(2)$ flavour fugacities recovers the
weak-Jacobi weight-$0$ index-$1$ form $\phi_{0,1}(q,y)$ carrying
$\chi(\mathrm{K3}) = 24$ at the elliptic-genus leading coefficient;
Borcherds singular-theta lifting of $\phi_{0,1}$ returns
$\Delta_5 = \mathrm{Grit}(\eta^9\vartheta_1)$. The two-step composite
arrow
\[
 \mathcal I_{\mathrm{Schur}}
 \;\xrightarrow{\;M_{24}\text{-avg over fugacity}\;}\;
 \phi_{0,1}(\mathrm{K3})
 \;\xrightarrow{\;\mathrm{Borch}\;}\;
 \Delta_5
\]
is the Gaiotto identification and is falsifiable at
Vol.~III Proposition~\ref{V3-prop:k3-schur-q-expansion}: any Fourier
coefficient mismatch through $q^{10}$ falsifies the
class-$\mathcal S$ origin.
\end{remark}

\section{Rank-\texorpdfstring{$N$}{N} K3 class-\texorpdfstring{$\mathcal S$}{S} family}
\label{sec:thqg-occ-supplement-ii}

\begin{remark}[The $N$-parametric family
 $\mathbf H^{(N)}_{\Delta_5}$ at $N \ge 2$]
\label{rem:thqg-occ-N-parametric}
Class-$\mathcal S$ is $N$-parametric: for each simply-laced Lie type
$\mathfrak g = A_{N-1}$ and each Riemann surface $C = \Sigma_{g,n}$,
Gaiotto~2009 (\texttt{arXiv:0904.2715}) constructs a $4d$
$\mathcal N = 2$ SCFT $\mathcal T[A_{N-1}, C]$ whose central charges
are computable from the Chacaltana--Distler pants-decomposition data
(Chacaltana--Distler 2010, \texttt{arXiv:1008.5203}, Sec.~5;
Chacaltana--Distler--Tachikawa 2013, \texttt{arXiv:1212.3952}). The
K3 base point $\mathcal T[A_1, \Sigma_{0,24}]$ extends to a
family
\[
 \bigl\{\mathbf H^{(N)}_{\Delta_5}
 := \text{Beem--Rastelli protected VOA of }
 \mathcal T[A_{N-1}, \Sigma_{0,24}]\bigr\}_{N \ge 2},
\]
indexed by $N$; the original $\mathcal T[A_1,\Sigma_{0,24}]$ is the $N = 2$ entry.
\end{remark}

\begin{proposition}[Chacaltana--Distler central charges for
 $\mathcal T\lbrack A_{N-1}, \Sigma_{0,24}\rbrack$]
\label{prop:thqg-occ-CD-ANm1-24}
\ClaimStatusProvedHere
For $g = 0$, $n = 24$ maximal regular $\mathfrak{su}(N)$ punctures,
the Gaiotto pants decomposition yields
$n_{\mathrm{trinions}} = 2g - 2 + n = 22$ and
$n_{\mathrm{tubes}} = 3g - 3 + n = 21$. The $T_N$ trinion anomaly
data (Chacaltana--Distler 2010 Tables~3,5; Benini--Tachikawa--Xie
2009 \texttt{arXiv:0906.0359}; Tachikawa 2013 lectures
\texttt{arXiv:1312.2684}~Ch.\ 13) interpolate to
\[
 n_v(T_N) = \frac{(N-1)(N-2)(N+2)}{2},
 \qquad
 n_h(T_N) = \frac{2N(N^2 - 1)}{3},
\]
verified at $N \in \{2, 3, 4, 5\}$ against the tabulated trinion
values $n_v = 0, 5, 18, 42$ and $n_h = 4, 16, 40, 80$ (CD 2010
Tab.\ 3; $N = 3$ is the Minahan--Nemeschansky $E_6$ SCFT with
$(a_{4d}, c_{4d}) = (41/24, 13/6)$, CD 2010 eq.\ (5.14)). Combining
with the $\mathrm{SU}(N)$ gauge vectors on the $n_{\mathrm{tubes}}$
tubes (each contributing $n_v = N^2 - 1$) gives the class-$\mathcal S$
count
\[
 n_v(A_{N-1}, \Sigma_{0,24}, \mathrm{all\ max})
 = 21(N^2 - 1) + 22\cdot\tfrac{(N-1)(N-2)(N+2)}{2},
\]
\[
 n_h(A_{N-1}, \Sigma_{0,24}, \mathrm{all\ max})
 = 22\cdot\tfrac{2N(N^2 - 1)}{3}.
\]
Applying the Shapere--Tachikawa normalisation (Shapere--Tachikawa
2008, \texttt{arXiv:0804.1957}) $a_{4d} = (5n_v + n_h)/24$,
$c_{4d} = (2n_v + n_h)/12$ and the Beem--Rastelli $4d/2d$
dictionary (BLLPRvR 2013, \texttt{arXiv:1312.5344}, eq.~(3.14))
$c_{2d} = -12\,c_{4d}$ yields
\begin{center}
\begin{tabular}{|c|c|c|c|c|c|c|}
\hline
$N$ & $n_v$ & $n_h$ & $a_{4d}$ & $c_{4d}$ & $c_{2d}$ &
 $\mathrm{rank}\,\mathcal M_{\mathrm{Coul}}$ \\
\hline
$2$ & $63$ & $88$ & $403/24$ & $107/6$ & $-214$ & $21$ \\
$3$ & $278$ & $352$ & $871/12$ & $227/3$ & $-908$ & $42$ \\
$4$ & $711$ & $880$ & $4435/24$ & $1151/6$ & $-2302$ & $63$ \\
\hline
\end{tabular}
\end{center}
with Coulomb rank $(N-1)(3g - 3 + n) = 21(N-1)$ and flavour
subalgebra $\mathfrak{su}(N)^{24}$ at Beem--Rastelli level
$k_{2d} = -N$ per puncture (BLLPRvR 2013 eq.~(3.18)).
\end{proposition}

\begin{proof}
For $A_1$ ($N = 2$), the boxed first-principles value
$c_{4d} = (5n - 13)/6$ at $n = 24$ returns $107/6$, cross-checked at
$n = 4$ against SU(2) $N_f = 4$ SQCD
$(a, c) = (23/24, 7/6)$ (Tachikawa~2013 lectures Ch.\ 13, eq.~(13.6))
and at $n = 3$ against the free trinion $T_2$ ($4$ free hypers,
$(a,c) = (1/6, 1/3)$). For $N = 3$ at $n = 3$, the formula returns
the $E_6$ Minahan--Nemeschansky values $(a, c) = (41/24, 13/6)$
(Minahan--Nemeschansky 1996, Argyres--Seiberg 2007
\texttt{arXiv:0711.0054}). The $N$-parametric polynomial
$(n_v(T_N), n_h(T_N))$ of
the Proposition's body is uniquely determined by the four-data
interpolation $(N, n_v, n_h)\in\{(2,0,4),(3,5,16),(4,18,40),(5,42,80)\}$;
higher-$N$ CD values (e.g.\ $T_6$) are consistent with the formula
(Chacaltana--Distler 2010 eq.\ (5.14)). Substitution into the
Shapere--Tachikawa formulas gives the table.
\end{proof}

\begin{remark}[$A_1$ central-charge normalization check]
\label{rem:thqg-occ-A1-central-charge-normalization}
The expression
$c_{4d}(A_1, \Sigma_{g,n}) = (12(g-1) + 7n)/6$
fails the $n = 4$ cross-check: at $(g, n) = (0, 4)$ it returns
$c_{4d} = 8/3$, whereas the well-established
$\mathrm{SU}(2)$ $N_f = 4$ value is $c_{4d} = 7/6$ (Shapere--Tachikawa
2008 Table~4; Tachikawa 2013 lectures eq.~(13.6)). The correct
first-principles value at $n = 24$ is therefore $c_{4d} = 107/6$
and $c_{2d} = -214$, coinciding with the boxed K3 value
in \texttt{k3\_chiral\_bialgebra\_platonic.tex}. The value
$c_{4d} = 26$ is incompatible with this normalization. The
alternative K3 value $a_{4d} = 421/24$ differs from the
first-principles $a_{4d} = 403/24$ by $18/24 = 3/4$; this belongs to
the $a$-anomaly normalization problem rather than the $c_{4d}$
calculation. The six-fold pentagon-umbral
$\mathbb{Z}/6$ anomaly is independent of this
central-charge correction (it lies in $H^3(M_{24}, U(1))$, the
generalised-Mathieu-moonshine 't~Hooft anomaly of
Gaberdiel--Persson--Ronellenfitsch--Volpato 2013, not in the
$c_{4d}$ data) and survives unchanged.
\end{remark}

\begin{remark}[$\mathfrak{su}(N)^{24}$ flavour subalgebra and
 $M_{24}$ permutation]
\label{rem:thqg-occ-flavour}
For each $N \ge 2$, the flavour subalgebra $\mathfrak{su}(N)^{24}
\subset \mathbf H^{(N)}_{\Delta_5}$ is an affine Kac--Moody
subalgebra at Beem--Rastelli level $k_{2d} = -N$ per
$\mathfrak{su}(N)$ factor (BLLPRvR 2013 Table~1, cross-verified
against the Minahan--Nemeschansky $E_6$ series
$k_{4d} = 6 \Rightarrow k_{2d} = -3$ at $N = 3$, C\'ordova--Shao
2015 \texttt{arXiv:1506.00265}). The Steiner $S(5,8,24)$ rigidity
of the $24$ punctures (Proposition~k3-steiner-rigidity in
Vol.~III) forces the puncture permutation symmetry to $M_{24}$
for every $N$, so the $\mathfrak{su}(N)^{24}$ subalgebra is
$M_{24}$-equivariant with the same projective anomaly cocycle of
order $6$ (Theorem~k3-projective-M24-equivariance, Vol.~III).
\end{remark}

\begin{remark}[Higher-rank Siegel lift: conjectural input data]
\label{rem:thqg-occ-siegel}
For $\mathcal T[A_1, \Sigma_{0,24}]$, the Borcherds singular-theta
lift of the $M_{24}$-averaged diagonal Schur index is the weight-$5$
Siegel cusp form $\Delta_5$ on $\mathrm{Sp}_4(\mathbb{Z})$ (paramodular
at level one; Gritsenko--Nikulin 1997
\texttt{arXiv:alg-geom/9504006}; Borcherds 1998
\texttt{arXiv:alg-geom/9609022} Theorem~13.3). For
$\mathcal T[A_{N-1}, \Sigma_{0,24}]$ at $N \ge 3$, no weight formula
is proved in this chapter. The conjectural automorphic lift must
specify four data before a weight can be asserted: the input Jacobi
or vector-valued modular form extracted from the $M_{24}$-averaged
Schur index, the Lorentzian lattice and paramodular group, the Weyl
vector, and the Borcherds constant term that gives the weight. The
$N=2$ case supplies the normalization $\Delta_5$ of weight~$5$;
higher-rank weights are independent of the rank formula
$\tfrac12(21(N-1)+2)$, which is already inconsistent with the $N=2$
normalization. The existence and automorphic identification of a
form $\Delta_{k_N}^{(N)}$ for $N \ge 3$ is conjectural. The
$N$-independent CY-to-chiral functor $\Phi$
(Vol.~III Chapter~\ref{ch:cy-to-chiral-functor}) predicts a
candidate source; the BKM-root structure at $N \ge 3$ requires data
beyond Mathieu moonshine and has no closed automorphic form supplied
in this chapter.
Primary references: Freitag ``Siegelsche Modulfunktionen'' 1983;
Gritsenko--Nikulin 1997--2018 series; Borcherds 1998 Theorems
13.3, 14.2 (general rank).
\end{remark}

\begin{remark}[K3-specificity of $(g, n) = (0, 24)$ at all $N$]
\label{rem:thqg-occ-k3-specificity}
The K3-specificity of the choice
$\mathcal T[A_1, \Sigma_{0,24}]$ resides entirely in the
puncture--surface datum $(g, n) = (0, 24)$: the number $24$
simultaneously realises
$\chi(\mathrm{K3}) = 24$ (Kodaira 1960),
the Euler number of the elliptic discriminant locus
$\#\{I_1\text{ fibres of generic elliptic K3}\} = 24$ (Kodaira 1963),
the length of the binary Golay code $\mathcal G_{24}$ and the Mathieu
$M_{24}$ dimension (Mathieu 1861; Curtis 1976; Conway--Sloane 1988),
and the Conway--Leech central charge $c = 24$
(Frenkel--Lepowsky--Meurman 1988). The Lie-type parameter $N$ is
\emph{independent} of this $(0, 24)$ specificity; for any $N \ge 2$
the construction yields an object in the K3-CY-family via the
CY-to-chiral functor
$\Phi(\mathrm{DCoh}(\mathrm{K3} \times E)) \to \mathbf H^{(N)}_{\Delta_5}$
(Vol.~III Chapter~\ref{ch:cy-to-chiral-functor}) with the rank-$N$
K3-CoHA as the $\Phi$-source. For other $(g, n)$ the Gaiotto
construction yields different objects (e.g.\
$\mathcal T[A_{N-1}, \Sigma_{g,n}]$ for $g \ge 1$) outside the
K3 branch selected by the $\Phi$-arrow, so the name
$\mathbf H^{(N)}_{\Delta_5}$ is
reserved for the $(g, n) = (0, 24)$ branch.
\end{remark}

\begin{remark}[CoHA-to-chiral comparison for
 $\mathbf H^{(N)}_{\Delta_5}$]
\label{rem:thqg-occ-coha-vs-chiral}
The rank-$N$ K3-CoHA $\mathrm{CoHA}_{\mathrm{K3}, N}$ (Kontsevich--Soibelman
2008 \texttt{arXiv:0811.2435}; Davison 2016
\texttt{arXiv:1603.00336}) is a \emph{cohomological Hall algebra}: an
associative $\mathbb{Z}$-graded algebra on
$\bigoplus_\gamma H^\star_T(\mathfrak{M}^{\mathrm{K3}}_{\gamma, N})$.
The chiral bialgebra $\mathbf H^{(N)}_{\Delta_5}$ is a different
object: an $E_1$-chiral bialgebra on the curve
$E^{\mathrm{nod,sm}}_{24}$ in the Ran topology. The CY-to-chiral
functor $\Phi$ (Vol.~III) supplies the comparison map
$\Phi(\mathrm{CoHA}_{\mathrm{K3},N}) \to
\mathbf H^{(N)}_{\Delta_5}$ on the Koszul locus; an equivalence is
the Vol.~III comparison theorem's hypothesis/conclusion rather than a
definition in this Vol.~I register. The rank-$N$ extension preserves
the distinction between a cohomological Hall algebra and an
$E_1$-chiral bialgebra.
\end{remark}
